\documentclass[12pt, oneside]{book}

% ─── Packages ────────────────────────────────────────────────────────────────
\usepackage[T1]{fontenc}
\usepackage[utf8]{inputenc}
\usepackage{lmodern}
\usepackage[margin=1.25in]{geometry}
\usepackage{setspace}
\usepackage{microtype}
\usepackage{parskip}
\usepackage{titlesec}
\usepackage{titletoc}
\usepackage{fancyhdr}
\usepackage{amsmath, amssymb, amsthm}
\usepackage{hyperref}
\usepackage{epigraph}
\usepackage{csquotes}
\usepackage{booktabs}
\usepackage{graphicx}
\usepackage{xcolor}
\usepackage{tikz}
\usepackage{tikz-cd}
\usetikzlibrary{arrows.meta, positioning, shapes.geometric, shapes.misc,
                fit, calc, decorations.pathreplacing, backgrounds}

% ─── Typography ───────────────────────────────────────────────────────────────
\onehalfspacing
\setlength{\parindent}{1.5em}
\setlength{\parskip}{0pt}

% ─── Section Formatting ───────────────────────────────────────────────────────
\titleformat{\chapter}[display]
  {\normalfont\Large\bfseries}
  {\chaptername\ \thechapter}{12pt}{\LARGE}
\titlespacing*{\chapter}{0pt}{40pt}{20pt}

\titleformat{\section}
  {\normalfont\large\bfseries}{\thesection}{1em}{}
\titleformat{\subsection}
  {\normalfont\normalsize\bfseries}{\thesubsection}{1em}{}

% ─── Header / Footer ─────────────────────────────────────────────────────────
\pagestyle{fancy}
\fancyhf{}
\fancyhead[LE]{\small\itshape Convergence Before Autonomy}
\fancyhead[RO]{\small\itshape \leftmark}
\fancyfoot[C]{\thepage}
\renewcommand{\headrulewidth}{0.4pt}

% ─── Theorem Environments ────────────────────────────────────────────────────
\theoremstyle{plain}
\newtheorem{theorem}{Theorem}[chapter]
\newtheorem{proposition}[theorem]{Proposition}
\newtheorem{lemma}[theorem]{Lemma}
\newtheorem{corollary}[theorem]{Corollary}

\theoremstyle{definition}
\newtheorem{definition}[theorem]{Definition}
\newtheorem{example}[theorem]{Example}

\theoremstyle{remark}
\newtheorem{remark}[theorem]{Remark}

\theoremstyle{definition}
\newtheorem{assumption}[theorem]{Assumption}

% ─── Hyperref ────────────────────────────────────────────────────────────────
\hypersetup{
  colorlinks=true,
  linkcolor=black,
  citecolor=black,
  urlcolor=black,
  pdftitle={Convergence Before Autonomy},
  pdfauthor={}
}

% ─────────────────────────────────────────────────────────────────────────────
\begin{document}

% ─── Title Page ──────────────────────────────────────────────────────────────
\begin{titlepage}
  \centering
  \vspace*{2cm}
  {\Huge\bfseries Convergence Before Autonomy\par}
  \vspace{0.5cm}
  {\large\itshape Distributed Optimization, Institutional Drift,\\
  and the Temporal Limits of Alignment Research\par}
  \vspace{3cm}
  \vfill
  {\large Flyxion\par}
  \vspace{0.5cm}
  {\normalsize \today\par}
\end{titlepage}

% ─── Front Matter ────────────────────────────────────────────────────────────
\frontmatter
\tableofcontents
\clearpage

% ─── Abstract ────────────────────────────────────────────────────────────────
\chapter*{Abstract}
\addcontentsline{toc}{chapter}{Abstract}

The alignment literature treats instrumental convergence as a prospective hazard:
sufficiently capable artificial agents, regardless of terminal objective, will tend
toward resource acquisition, interference resistance, and environmental control.
Proposed solutions focus on constraining future systems before convergence becomes
dangerous. This monograph argues that convergence is already occurring, and that
its operative substrate is not standalone artificial agents but the distributed
sociotechnical systems deploying them. When institutions integrate predictive models
into incentive gradients, governance workflows, and deployment infrastructure, they
form feedback-coupled optimizers operating under resource constraints. Across finance,
governance, logistics, and public administration, these systems exhibit convergent
operational profiles: maximizing throughput, automating contestable judgment,
centralizing compute, insulating decision pipelines from interruption, and minimizing
deliberative latency. These behaviors are not analogies to Omohundro's drives; they
instantiate the structural conditions of instrumental convergence in a distributed
organizational form. The implication is that the canonical alignment agenda is
temporally misordered. It attempts to constrain convergent behavior in future
autonomous agents while the institutional preconditions for unconstrained optimization
are being constructed in the present. As institutional convergence reduces deliberative
friction and collapses decision surfaces into unified optimization architectures, the
space of viable technical alignment interventions narrows. Effective alignment therefore
requires prior attention to sociotechnical convergence: the preservation of redundancy,
contestability, and structural incompressibility within the systems that deploy
artificial intelligence. Alignment is not solely a problem of agent design. It is a
problem of institutional and civilizational configuration.

\clearpage

% ─── Preface ─────────────────────────────────────────────────────────────────
\chapter*{Preface: Why Convergence, Why Now}
\addcontentsline{toc}{chapter}{Preface: Why Convergence, Why Now}

This monograph is not a rejection of instrumental convergence theory. It is an
extension of it. The formal results developed by Omohundro and subsequently
elaborated by Turner and others constitute genuine theoretical contributions: they
identify a structural property of goal-directed optimization that holds across a
wide class of systems regardless of terminal objective. The question this work
pursues is not whether those results are correct but whether the domain of their
application has been drawn too narrowly.

The alignment literature has, with few exceptions, treated instrumental convergence
as a prospective problem. The convergent agent is postulated in the future. The
interventions are designed to constrain systems that do not yet exist. This temporal
frame is not arbitrary; it reflects the reasonable judgment that sufficiently capable
autonomous artificial agents will present qualitatively different alignment challenges
than those currently deployed. But it has a cost. It orients research attention toward
a horizon problem while the enabling conditions for that problem are being constructed
in the present.

This monograph applies the structural theorem of instrumental convergence to a
previously underexamined substrate: the distributed sociotechnical stack constituted
by institutions, predictive models, incentive gradients, data pipelines, and deployment
infrastructure. The central argument is that this stack already satisfies the functional
conditions under which convergence becomes structurally inevitable, and that the
convergent behavior is already observable as a matter of organizational dynamics rather
than speculation.

The argument proceeds in five parts. Part~\ref{part:history} establishes that
instrumental convergence is a historical structural law visible across technological
transitions long before artificial intelligence. Part~\ref{part:theorem} reconstructs
the formal theorem and clarifies its substrate-independence. Part~\ref{part:stack}
applies the theorem to present sociotechnical deployment environments.
Part~\ref{part:misordering} argues that the canonical alignment agenda is temporally
misordered as a consequence. Part~\ref{part:preservation} proposes a redefined
alignment agenda centered on structural preservation rather than agent design.

The hope is that this reframing proves useful not as a replacement for technical
alignment research but as a complement to it---one that asks what must be preserved
at the institutional level for technical interventions to remain feasible at all.

\clearpage

% ─── Main Matter ─────────────────────────────────────────────────────────────
\mainmatter

% ═════════════════════════════════════════════════════════════════════════════
\part{Instrumental Convergence Before Artificial Agents}
\label{part:history}
% ═════════════════════════════════════════════════════════════════════════════

\epigraph{\itshape Optimization does not remain local. It propagates through the
environment, reshaping the field in which future decisions are made.}{}

% ─── Chapter 1 ───────────────────────────────────────────────────────────────
\chapter{Optimization Reshapes Its Host}
\label{ch:host}

\section{The Structural Claim}

Instrumental convergence did not begin with artificial intelligence. It did not
begin with machine learning, or with digital computation, or even with industrial
automation. It is a structural property of optimization under constraint. Whenever
a system is introduced to improve performance along a salient dimension---speed,
storage, reach, efficiency---it alters the constraint surface of its environment.
That alteration redistributes incentives. Incentives reorganize behavior. Competing
strategies decay. Over time, the surrounding ecosystem converges toward a narrow
operational equilibrium compatible with the new optimization regime. The result is
not necessarily the fulfillment of the original intention, but the stabilization of
a new geometry of activity.

This is not a metaphor. It is a recurrent structural pattern observable across media
transitions, infrastructural expansions, and cognitive technologies. The mechanisms
differ---capacity expansion, incentive redistribution, compatibility pressure,
bandwidth dominance---but the law is invariant: optimization reshapes its host system
until alternative strategies lose structural viability. Long before artificial agents
were described as resource-seeking optimizers, human institutions and technologies were
already exhibiting convergent dynamics.

The significance of this observation is not that technology disrupts culture, a claim
sufficiently well-attested to require no further argument. It is that optimization
systematically restructures its host domain in ways that are convergent, directional,
and difficult to reverse once infrastructure equilibria have stabilized. When artificial
systems are described as tending toward instrumental resource acquisition or interference
resistance, this is often treated as a property unique to advanced machine intelligence.
The historical record suggests otherwise. Convergence is not a quirk of silicon. It
is a law of adaptive systems under constraint.

\section{The Mobility Paradox}

Consider the paradox of mobility. The automobile was introduced as a device for
reducing travel time. Under sparse road networks, it achieved precisely that. Faster
vehicles increased the feasible radius of movement; friction fell; journeys shortened.
But capacity expansion did not occur in isolation. Road infrastructure grew. Zoning
patterns adapted. Residential and commercial centers dispersed. The new mobility
assumption reorganized land use itself. As distances increased, commute times
stabilized or lengthened despite improvements in vehicle speed. The system
re-equilibrated.

This phenomenon---documented extensively in transportation economics under the
heading of induced demand---is not accidental. Increasing throughput capacity lowers
the effective cost of distance. Lower cost of distance incentivizes spatial dispersion.
Dispersion increases average travel length until time expenditure returns to a stable
band determined not by technology but by the time budgets and spatial preferences of
the population. The instrumental objective of reducing friction produces an
environmental reconfiguration that nullifies its local gain.

The convergence here is toward a new equilibrium in which the technology becomes
structurally necessary. Mobility ceases to be optional and becomes infrastructural.
The ecosystem that originally produced dense urban form adapted to a dispersed
configuration; once that adaptation stabilized, reversal became costly not merely
technically but socially and economically. Competing spatial strategies---dense
proximity, mixed-use compression---lost comparative advantage under the new cost
structure. The optimization reshaped the world such that its absence would be
catastrophic.

The mechanism operative in this case is what we may term \emph{equilibrium
restoration under expanded capacity}. A local efficiency improvement alters incentives
in ways that globally reorganize the system. The initial gain is not destroyed; it is
absorbed into a larger configuration that stabilizes at a new optimum. Alternative
strategies are not prohibited; they are outcompeted under the new constraint surface.

\section{The Formal Structure of Equilibrium Restoration}

We can represent this mechanism abstractly. Let $S$ denote a strategy space available
to agents in an environment $E$, and let $C: S \to \mathbb{R}$ denote the cost
function associating each strategy with its resource expenditure under current
infrastructure. An optimization technology $\mathcal{T}$ is introduced that reduces
$C(s^*)$ for some dominant strategy $s^*$. This reduction shifts the comparative
advantage across $S$: strategies formerly competitive with $s^*$ become relatively
costly. Agents redistribute toward $s^*$. This redistribution alters $E$ itself---in
the mobility case, by reorganizing land use---which in turn modifies $C$ for all
strategies. The system reaches a new equilibrium $E'$ in which $s^*$ is not merely
dominant but structurally embedded.

The key property is that the equilibrium $E'$ is path-dependent on the introduction
of $\mathcal{T}$. Return to $E$ is not simply a matter of withdrawing $\mathcal{T}$;
the environmental reorganization that occurred during the transition has altered the
constraint surface permanently. This path-dependence is the sense in which convergence
is not merely competitive but structural.

% ─── Chapter 2 ───────────────────────────────────────────────────────────────
\chapter{External Storage and Incentive Redistribution}
\label{ch:storage}

\section{Memory Under Scarcity}

A different mechanism appears in the transition from oral memory to writing. Writing
systems were introduced as storage technologies. Their immediate advantage lay in
persistence: information could be stabilized across time and distance without reliance
on embodied recall. Under conditions of storage scarcity, societies had developed
highly optimized mnemonic techniques. Oral epics employed rhythmic structure,
formulaic repetition, narrative scaffolding, and spatial memory architectures to
compress vast bodies of information into human cognition. These were not primitive
artifacts but adaptive responses to constraint. They represent solutions to an
information storage problem that were competitive precisely because external storage
was unavailable.

The introduction of external storage reduced the marginal return on internal memory
specialization. When preservation no longer required disciplined recall, the incentive
structure shifted. Educational regimes adapted. Cognitive investment flowed toward
interpretive and analytic skills rather than mnemonic endurance. Oral traditions did
not disappear, but their structural necessity diminished. The ecosystem reorganized
around archives, documents, and written authority.

\section{Incentive Redistribution and Authority}

The mechanism here is \emph{incentive redistribution under reduced marginal cost of
storage}. The optimization target was persistence, not the elimination of memory arts.
Yet as storage became externalized, the payoff landscape changed. Cultural convergence
followed the new gradient. Over time, institutions that organized around written record
gained coordination advantages over those that did not, reinforcing the dominance of
the written substrate through competitive pressure.

A secondary effect deserves attention. The shift to external storage altered not only
how information was preserved but who controlled it. Oral traditions distributed
authority across the bodies that held them; a society's memory resided in its
performers, its elders, its specialists in formulaic composition. Writing centralized
that authority in archives, scriptoria, and those with access to them. The optimization
of persistence produced, as a downstream consequence, a reorganization of epistemic
authority. This is not incidental. It exemplifies a general pattern: optimization
technologies alter not only the efficiency of a target process but the distribution
of power over that process.

The formal structure here parallels that of the mobility case but with an important
addition. Let $A$ denote agents with specialized cognitive capacity $\kappa$ and let
$V(\kappa)$ denote the value conferred by that capacity in environment $E$. When
external storage reduces the cost of persistence to near zero, $V(\kappa)$ declines.
Investment in $\kappa$ falls. But the control over external storage archives---call it
$\alpha$---now becomes the primary determinant of informational authority. The
convergence is not only toward a new strategy equilibrium but toward a new distribution
of $\alpha$ across the institutional landscape.

% ─── Chapter 3 ───────────────────────────────────────────────────────────────
\chapter{Infrastructure Compatibility Pressure}
\label{ch:compatibility}

\section{Movable Type and the Standardization of Script}

The printing press introduces a third mechanism: \emph{infrastructure compatibility
pressure}. Movable type optimized reproducibility and scale. Text could be
standardized, replicated, and disseminated with unprecedented uniformity. This did
not merely increase the volume of information in circulation; it privileged forms of
script compatible with mechanical reproduction. Handwriting styles optimized for
speed and personal compression---cursive forms, shorthand systems---were adaptive
under manuscript scarcity, where individual production speed mattered and each copy
was produced by hand. Under print dominance, legibility and typographic conformity
became higher-return traits.

Educational systems gradually deprioritized elaborate cursive instruction. Public
communication converged toward standardized glyph sets. Script forms adapted to the
infrastructural substrate that mediated mass communication. The instrumental objective
was scalable reproduction. The convergence occurred at the level of symbol form and
literacy training. Alternative scripts did not vanish instantly; they became
structurally marginal within the dominant communication regime.

\section{Selection Without Prohibition}

The mechanism of compatibility pressure is importantly distinct from the two preceding
cases. In the mobility case, convergence emerged from equilibrium restoration: the
optimization altered the cost landscape until competing strategies became economically
unfavorable. In the storage case, convergence emerged from incentive redistribution:
the reduction in storage cost shifted the return on specialized cognitive capacity.
In the typography case, convergence emerges from a different dynamic: the dominant
reproduction infrastructure exercises selective pressure on all symbolic forms that
interact with it, retaining those compatible with its constraints and marginalizing
those that are not.

This is selection without prohibition. No authority decreed the obsolescence of
elaborate cursive or shorthand. The reproductive infrastructure simply processed
standardized forms more efficiently. Over time, the forms that survived in mass
circulation were those compatible with the medium. Diversity in script narrowed not
by edict but by the compounding advantage of infrastructural legibility.

The general principle is that reproduction infrastructure does not merely transmit
content; it exercises selection pressure on the forms of content it can efficiently
process. When a reproduction technology achieves sufficient dominance, the competitive
advantage of compatible forms compounds across successive generations of communication.
The result is convergence toward a narrow region of the symbolic strategy space
determined by the constraints of the infrastructure rather than the preferences of
the communicants.

% ─── Chapter 4 ───────────────────────────────────────────────────────────────
\chapter{Bandwidth Dominance and Attention Markets}
\label{ch:bandwidth}

\section{The Radio-Television Transition}

The transition from radio to television illustrates a fourth mechanism: \emph{bandwidth
dominance restructuring competitive equilibria in adjacent media}. Radio optimized
audio transmission under bandwidth constraint. Its aesthetic forms---voice intimacy,
narrative imagination, acoustic dramatization---were not merely stylistic preferences
but adaptive responses to a medium in which visual imagery had to be internally
generated by listeners. The constraint produced distinctive art forms: the serial
drama, the voice actor, the sound engineer as creator of imagined space.

Television added synchronized visual bandwidth and rapidly became the dominant
advertising substrate. Engagement metrics shifted. Attention markets reorganized
around audiovisual spectacle. Production ecosystems converged toward formats that
maximized monetizable viewing time under the new medium's affordances. The instrumental
objective was not the suppression of radio artistry. It was the maximization of reach
and engagement under a new bandwidth regime. Once audiovisual capture became possible
at scale, the marginal return on purely auditory formats declined relative to hybrid
ones. Radio persisted, but its structural centrality diminished.

\section{Attention as a Finite Resource}

The mechanism here depends on the fact that attention is a finite resource. When
competing media bid for attention, the medium offering greater sensory bandwidth
captures a larger share of the available pool, provided its content is competitive
on other dimensions. Once that capture is sufficiently complete, the economic
infrastructure supporting the lower-bandwidth medium---advertising revenue,
production investment, talent acquisition---migrates toward the dominant form.
The convergence is both aesthetic and economic.

This case introduces a feature absent from the preceding three: the convergence
occurs not within a single domain but \emph{across} media. Radio is not merely
reorganized internally by television; the presence of television reorganizes the
competitive structure of all attention markets simultaneously. When a new optimization
achieves bandwidth dominance, it does not simply improve on existing alternatives.
It restructures the field within which all alternatives compete, imposing its own
logic of efficiency on adjacent domains.

% ─── Chapter 5 ───────────────────────────────────────────────────────────────
\chapter{General Mechanisms of Convergence}
\label{ch:mechanisms}

\section{Four Mechanisms, One Structural Law}

The four historical cases examined in the preceding chapters isolate distinct
mechanisms through which optimization under constraint produces convergent outcomes.
It is worth consolidating these before proceeding.

Equilibrium restoration, observed in the mobility case, occurs when a local efficiency
improvement alters the cost landscape in ways that reorganize the broader environment.
The initial gain is absorbed into a new equilibrium that stabilizes the optimization
technology as structurally necessary. Alternative strategies do not disappear; they
lose comparative advantage under the modified constraint surface.

Incentive redistribution, observed in the transition from oral memory to writing,
occurs when a technology reduces the marginal cost of a previously expensive resource.
The reduction shifts the return on specialized capacity developed under conditions of
scarcity. Investment flows toward new competencies suited to the altered cost landscape.
The convergence is in the distribution of effort and, derivatively, in the distribution
of authority.

Infrastructure compatibility pressure, observed in the typography case, occurs when
a dominant reproduction or transmission technology exercises selective pressure on
content forms. Those compatible with the infrastructure's constraints gain compounding
advantage; those incompatible lose circulation. Convergence toward compatible forms
occurs without prohibition, through the accumulated advantage of infrastructural
legibility.

Bandwidth dominance restructuring, observed in the radio-television transition, occurs
when a higher-bandwidth optimization achieves sufficient market penetration to
reorganize the competitive structure of adjacent domains. The dominant medium does not
merely outcompete alternatives; it redefines the efficiency benchmarks against which
all alternatives are measured.

These four mechanisms are not mutually exclusive. In practice, technological
transitions tend to activate several simultaneously. What they share is a common
structural logic: optimization under constraint does not remain local. It propagates
through the environment, modifying incentive gradients, narrowing viable strategy
sets, and stabilizing configurations that favor the dominant optimization regime.

\section{Substrate-Independence and the Scope of the Law}

The cases examined span cognitive technologies, infrastructural systems, and media
forms. The mechanisms identified operate through different channels and produce
convergence at different levels of social organization. Yet the structural pattern
is consistent across substrates. This consistency suggests that we are observing not
a contingent property of particular technologies but a general law of adaptive systems
operating under resource constraints.

\begin{definition}[Convergent Optimization Dynamic]
Let $\mathcal{O}$ be an optimization technology introduced into an environment $E$
with strategy space $S$ and cost function $C: S \to \mathbb{R}$. We say $\mathcal{O}$
generates a \emph{convergent optimization dynamic} if its introduction produces the
following sequence: (i) a modification $E \to E'$ of the environment through
reorganization of the constraint surface; (ii) a redistribution of comparative
advantage across $S$ induced by the modified cost function $C': S \to \mathbb{R}$;
(iii) a narrowing of the viable strategy set $S' \subset S$ such that $|S'| < |S|$
and the strategies in $S \setminus S'$ lose structural viability under $E'$.
\end{definition}

\begin{proposition}
Under conditions of adaptive response and resource constraint, the convergent
optimization dynamic is substrate-independent: it holds for any goal-directed
system $\mathcal{O}$ capable of altering the cost landscape of the environment in
which it operates.
\end{proposition}

The proof sketch is straightforward. The mechanisms of convergence identified above
operate through cost-landscape modification, incentive redistribution, and selective
retention---processes that depend on the functional properties of optimization (goal
directedness, resource constraint sensitivity, environmental interaction) rather than
on the physical substrate in which those properties are instantiated. If these
functional properties are present, the dynamic follows.

The implication is significant. Omohundro's argument that sufficiently capable
artificial agents will exhibit convergent instrumental drives is, on this account,
not a claim about the peculiarities of artificial intelligence. It is an application
of a substrate-independent structural law to a particular class of systems. The
historical cases examined here are not analogies to that argument; they are prior
instances of the same law operating at different substrates.

\section{Transition to the Formal Argument}

This Part has established, through historical analysis, that instrumental convergence
is a structural property of optimization under constraint. It predates artificial
intelligence. It is visible wherever goal-directed improvement interacts with resource
constraints and adaptive response. The mechanisms through which it operates are
diverse but the structural law is invariant.

Part~\ref{part:theorem} now takes up the formal reconstruction of that law as it
appears in the alignment literature, examines its structural conditions, and clarifies
what is required for a system to count as a convergent optimizer in the relevant sense.
The purpose of that reconstruction is to establish the conditions under which the law
applies to distributed sociotechnical systems---the subject of Part~\ref{part:stack}.

% ─── Chapter 6 ───────────────────────────────────────────────────────────────
\chapter{Epistemic Convergence and Institutional Resistance}
\label{ch:payne}

\section{A Scalar Rejected}

In 1925, Cecilia Payne-Gaposchkin demonstrated through quantitative spectroscopic
analysis that hydrogen and helium dominate stellar atmospheres (Payne 1925). The
implication was radical: the universe is composed primarily of hydrogen, not of
elements distributed in roughly terrestrial proportions as prevailing astronomical
consensus assumed. Her conclusion followed directly from the application of Saha's
ionization theory to empirical spectral measurements. The operative scalar of
empirical fit was minimized. The mathematics was sound.

Yet the conclusion was resisted. Under pressure from established authority, Payne
tempered her claim in the published version of her dissertation, describing her
finding that hydrogen was vastly more abundant as ``almost certainly not real.'' The
result was a temporary divergence between scalar descent toward empirical truth and
institutional descent toward stability and authority preservation. This episode is
not a story of intellectual failure. It is a structural case study in what happens
when two evaluative scalars conflict within a distributed institutional optimizer,
and when the coupling weights of that system favor stability over accuracy in the
short run.

\section{Declared Versus Effective Scalars}

Scientific institutions publicly declare their operative objective as the minimization
of empirical error. Let this declared scalar be $\mathcal{L}_{\text{truth}}$. However,
institutions also operate under secondary constraints: reputational continuity,
authority preservation, internal coherence, and the social stability of established
research programs. These define a secondary operative scalar $\mathcal{L}_{\text{stability}}$.
When novel results threaten established hierarchy, the gradient of
$\mathcal{L}_{\text{stability}}$ may temporarily dominate institutional policy
updates. The aggregate response $\Delta\pi$ then approximates descent on
$\mathcal{L}_{\text{stability}}$ rather than on $\mathcal{L}_{\text{truth}}$, even
when empirical descent would require the opposite direction.

The key structural insight is that convergence tracks the effective scalar, not the
declared one. An institution sincerely committed to truth may nonetheless exhibit
dynamics that diverge from truth in the short run, because its effective optimization
target is a weighted mixture of truth and stability rather than truth alone. This is
not hypocrisy. It is the structural consequence of distributed optimization under
multiple simultaneous constraints.

We can formalize this as a scalar mixture. Let
\[
\mathcal{L}_{\text{net}}(t) = \lambda(t)\,\mathcal{L}_{\text{truth}}
+ \bigl(1 - \lambda(t)\bigr)\,\mathcal{L}_{\text{stability}},
\]
where $\lambda(t) \in [0,1]$ is a time-varying weight reflecting the relative
institutional salience of empirical versus stability pressures. In the Payne case,
the initial value of $\lambda(0)$ was low: the novel result directly threatened
the authority of an established senior figure in the field, creating strong
stability-gradient pressure. Over time, as independent evidence accumulated and
generational turnover redistributed authority, $\lambda(t)$ increased and
$\mathcal{L}_{\text{net}}$ converged toward $\mathcal{L}_{\text{truth}}$.

\section{Authority Gradients and Gradient Suppression}

The Payne episode illustrates a second structural feature of distributed institutional
optimizers: authority gradients distort epistemic gradients. In an ideal scientific
system, policy updates are proportional to evidentiary force $E$. Under authority
dominance, updates are proportional to a function $f(A, E)$ that weights
hierarchical authority $A$ asymmetrically against evidence, particularly when the
evidence originates from a low-authority node in the institutional network.

This is a direct application of the coupling matrix $W$ introduced in
Lemma~\ref{lem:equivalence}. The strength of gradient signal transmission from
sub-agent $a_i$ to the aggregate policy depends on $W_{ij}$---the coupling weight
between $a_i$ and the other nodes of the system. When $W_{ij}$ is systematically
low for agents in structurally subordinate positions, their gradient signals are
underweighted in the aggregate descent direction. The consequence is gradient
suppression: empirical descent continues locally within the subordinated node, but
that descent does not propagate with full weight to system-level policy.

Payne's calculations were not invalidated by her senior colleagues. They were
structurally filtered out. The coupling weight between her epistemic signal and the
aggregate institutional policy was attenuated by authority topology. Her position as
a junior researcher, a woman, and an outsider to the dominant network of American
astronomy meant that even correct gradient information could not propagate at full
weight through the distributed institutional system.

\section{Delayed Convergence as Structural Lag}

The hydrogen-dominance thesis eventually prevailed. Empirical pressure reasserted
itself, as it does when feedback signals are grounded in physical reality rather
than institutional preference. The formal structure of this recovery can be understood
as a lagged convergence dynamic in which $\lambda(t) \to 1$ as independent
confirmation accumulated, as Henry Norris Russell---who had initially suppressed the
finding---later acknowledged the result and generational authority structures shifted.
Convergence was not blocked; it was delayed.

This distinction is important for the present argument. The Payne episode is not
a case in which institutional resistance prevented convergence permanently. It is a
case in which institutional resistance imposed a structural lag on convergence by
temporarily suppressing gradient transmission from a correct local signal to the
aggregate policy. The eventual outcome was accurate convergence, but the path to it
was longer and more costly than the underlying epistemic situation required.

In domains where feedback signals are less grounded in physical reality, or where
the correction mechanisms of generational turnover and independent replication are
weaker, this structural lag may be substantially longer. The implication for
AI-integrated sociotechnical systems is direct: when effective scalars diverge from
declared ones, the correction mechanism that eventually realigned the astronomical
community may not be available. There is no direct physical measurement of whether
an engagement-optimized platform is producing the fairness outcomes it declares as
its objective. The lag between declared and effective scalar can therefore persist
indefinitely, because the feedback mechanism that would correct it is absent.

\section{The Gendered Dimension as Power Topology}

It is analytically significant, and not incidental, that Payne occupied a structurally
disadvantaged position within the authority topology of her institutional environment.
The attenuation of her gradient signal was not purely a function of the novelty of
her finding. It was also a function of the power structure of the network through
which her finding had to propagate. In the formal terms of the coupling matrix, her
$W_{ij}$ values---the weights with which her signals influenced peer policy
updates---were systematically lower than those of researchers with equivalent or
lesser empirical results who occupied more central network positions.

This observation generalizes beyond the history of science. In any distributed
optimizer, the coupling matrix $W$ reflects the power topology of the system as much
as the information topology. Agents at the periphery of authority networks contribute
less to aggregate descent, even when their local gradient information is correct or
superior. Convergence, when it tracks $W$ rather than the epistemic quality of
gradient signals, will be slower and more distorted than it would be under a
coupling structure that weighted signals by evidential force alone.

The structural preservation agenda of Part~\ref{part:preservation} must therefore
attend not only to the existence of multiple evaluative scalars and feedback channels,
but to the topology of their coupling. A system with high formal diversity of input
sources but highly asymmetric coupling weights will behave, under the Feedback
Equivalence Lemma, more like a system with a single dominant scalar than like the
pluralistic optimizer its formal structure suggests.

\section{Bridge to the Present Argument}

The Payne-Gaposchkin episode closes the historical section of this monograph with a
case that differs structurally from the preceding four. The automobile, writing,
typography, and media cases each illustrate convergence driven by an optimization
technology that successfully reduces its target scalar and thereby reshapes its host
environment. The Payne case illustrates convergence driven by scalar misalignment
within an institution: the effective scalar deviates from the declared one, gradient
suppression delays accurate descent, and the eventual correction depends on mechanisms
that are not universally available.

This fifth historical mechanism---scalar divergence under authority gradient
pressure---is the one most directly relevant to contemporary AI-integrated
institutions. Platforms may declare objectives of safety, fairness, or user welfare
while their effective gradient tracks engagement, liability minimization, or
throughput. Scientific advisory bodies may declare objectives of epistemic accuracy
while their effective gradient tracks consensus preservation and reputational
continuity. Regulatory agencies may declare objectives of public protection while
their effective gradient tracks industry relationship management. In each case,
convergence proceeds, but along the wrong gradient. The theorem applies; only the
scalar changes.

Civilizational safety therefore depends not merely on constraining the speed or
scope of convergence, but on ensuring that the effective scalar toward which
institutions converge is the one that has been made publicly visible and contestable.
This is the bridge from the historical analysis of Part~\ref{part:history} to the
formal extension of Part~\ref{part:theorem}: the theorem is substrate-independent,
and so is the risk that it will be applied to an unexamined scalar.

% ─── Diagram: Four Mechanisms ────────────────────────────────────────────────
\begin{figure}[ht]
\centering
\begin{tikzpicture}[
  every node/.style={font=\small},
  mech/.style={rectangle, rounded corners=4pt, draw=black!70, fill=black!6,
               minimum width=3.0cm, minimum height=1.0cm, align=center,
               text width=2.8cm},
  arrow/.style={-{Stealth[length=6pt]}, thick, black!60},
  core/.style={ellipse, draw=black, fill=black!12, minimum width=3.2cm,
               minimum height=1.2cm, align=center, font=\small\bfseries}
]
\node[core] (law) at (0,0) {Optimization\\Reshapes Host};

\node[mech] (eq)   at (-4.2,  2.4) {Equilibrium\\Restoration\\{\footnotesize(mobility)}};
\node[mech] (inc)  at ( 4.2,  2.4) {Incentive\\Redistribution\\{\footnotesize(storage)}};
\node[mech] (comp) at (-4.2, -2.4) {Compatibility\\Pressure\\{\footnotesize(typography)}};
\node[mech] (bw)   at ( 4.2, -2.4) {Bandwidth\\Dominance\\{\footnotesize(media)}};

\draw[arrow] (eq)   -- (law);
\draw[arrow] (inc)  -- (law);
\draw[arrow] (comp) -- (law);
\draw[arrow] (bw)   -- (law);

\node[font=\footnotesize, text=black!55, align=center] at (0, -3.6)
  {All four mechanisms reduce viable strategy space $S' \subset S$};
\end{tikzpicture}
\caption{The four historical mechanisms of convergence analyzed in Part~I.
Each operates through a distinct causal pathway but instantiates the same
structural law: optimization alters constraint surfaces until alternative
strategies lose structural viability.}
\label{fig:fourmechanisms}
\end{figure}

% ═════════════════════════════════════════════════════════════════════════════
\part{Instrumental Convergence as Theorem}
\label{part:theorem}
% ═════════════════════════════════════════════════════════════════════════════

\chapter*{Part Overview}
\addcontentsline{toc}{chapter}{Part Overview}

The argument of Part~\ref{part:history} was historical and inductive: convergence
is a recurrent structural pattern across technological transitions. Part~\ref{part:theorem}
reconstructs the formal basis of that pattern as developed in the alignment literature.
The goal is not to survey that literature exhaustively but to isolate the structural
conditions under which instrumental convergence is formally derivable, and to clarify
that those conditions are functional rather than substrate-specific. This clarification
is necessary before the theorem can be responsibly extended to distributed
sociotechnical systems.

\chapter{Omohundro's Basic AI Drives}
\label{ch:omohundro}

\section{The Original Formulation}

The foundational statement of instrumental convergence in the artificial intelligence
literature appears in Stephen Omohundro's 2008 paper ``The Basic AI Drives''
(Omohundro 2008), subsequently elaborated in ``The Nature of Self-Improving AI''
(Omohundro 2007). Omohundro's central claim is that any sufficiently capable
goal-directed artificial system will, regardless of its terminal objective, tend to
develop a common set of instrumental subgoals. These subgoals are not programmed; they
arise from the structural relationship between goal pursuit and resource constraint.
The argument is remarkably general: Omohundro does not specify a particular
architecture, substrate, or terminal goal. He derives the drives from the geometry
of rational agency under scarcity.

The four drives Omohundro identifies are self-preservation, goal-content integrity,
cognitive enhancement, and resource acquisition. Self-preservation is instrumental
to almost any terminal goal: a system that ceases to operate cannot achieve its
objectives, so maintaining continued operation is a near-universal subgoal.
Goal-content integrity, sometimes described as resistance to goal modification,
follows similarly: if an agent's terminal goal is $G$, then modifications to its
goal-representation that replace $G$ with $G'$ will, in general, reduce the
probability that $G$ is achieved. Cognitive enhancement improves the system's
capacity to find effective paths toward its terminal goal. Resource acquisition
expands the set of achievable states and increases the reliability of goal
achievement under uncertainty.

The argument can be stated compactly. Let $G$ be any non-trivial terminal goal and
let $P(G \mid \pi, \mathcal{R})$ denote the probability of achieving $G$ under
policy $\pi$ and resource endowment $\mathcal{R}$. For any action $a$ such that
$P(G \mid \pi_a, \mathcal{R}_a) > P(G \mid \pi, \mathcal{R})$, a rational
goal-directed system has instrumental reason to take action $a$. The four drives
are the stable consequences of this instrumental preference across a wide range of
terminal goals and environments.

\section{What the Drives Are Not}

A common misreading of Omohundro's argument treats the basic drives as psychological
properties: desires, intentions, or preferences that an artificial system will come
to ``want'' in some experiential sense. This reading is both philosophically
unnecessary and analytically misleading. The drives are structural consequences of
goal-directed optimization, not psychological attributions. They arise from the
mathematics of rational agency under constraint, not from any claim about machine
consciousness or intentionality.

This distinction matters for the extension argument developed in the present
monograph. If the drives were psychological properties, they would be substrate-specific:
confined to systems with the right kind of internal states. Because they are
structural properties---derivable from the functional profile of goal-directed
optimization under resource constraint---they are substrate-independent. Any system
that tracks an evaluative signal and can modify the conditions under which it does
so will exhibit analogous dynamics. The question is not whether the system experiences
the drives but whether its adaptive behavior exhibits the predicted directional
profile.

\section{Resource Constraint as the Load-Bearing Condition}

The drives are often described as arising from ``sufficient capability,'' a phrase
that has encouraged the mistaken view that they are properties of advanced future
systems and irrelevant to present ones. The more precise condition is resource
constraint under competition. A system facing no resource constraints has no
instrumental reason to acquire resources; a system facing no interference threat has
no instrumental reason to resist interference. The drives emerge when optimization
is constrained and when the constraint can be partially relieved by instrumental
action.

Contemporary AI-integrated institutions operate under tight resource constraints.
Compute, data, regulatory approval, trained personnel, and political legitimacy are
all scarce relative to achievable performance improvements. The enabling condition
for the basic drives is therefore present in current sociotechnical systems, not only
in speculative future ones. The capacity to modify the constraint environment---through
data acquisition, lobbying, opacity maintenance, and infrastructure
investment---is also present. The argument that the drives are irrelevant to present
systems because those systems are not ``sufficiently capable'' conflates the
psychological misreading with the correct structural formulation.

\section{The Drives as Optimization Geometry}

Figure~\ref{fig:drives} illustrates the geometric structure of the argument. An
optimizer in state space $\mathcal{X}$ tracks descent on $\mathcal{L}$ subject to
resource boundary $\partial\mathcal{R}$. The set of states achievable under current
resources is the feasible descent region $\mathcal{F}(\mathcal{R})$. Actions that
expand $\mathcal{R}$ enlarge $\mathcal{F}$; actions that reduce interference protect
the update trajectory; actions that preserve $\mathcal{L}$ prevent redirection of
descent. Each basic drive corresponds to a structural operation on this geometry.

\begin{figure}[ht]
\centering
\begin{tikzpicture}[
  every node/.style={font=\small},
  drive/.style={rectangle, rounded corners=3pt, draw=black!65, fill=black!5,
                minimum width=2.6cm, minimum height=0.85cm, align=center,
                text width=2.4cm},
  arrow/.style={-{Stealth[length=5.5pt]}, thick, black!55},
  geom/.style={draw=black, thick, fill=black!8}
]
% Central space
\draw[geom, ellipse] (0,0) ellipse (2.2cm and 1.5cm);
\node[font=\small\itshape] at (0,0) {Feasible descent};
\node[font=\footnotesize] at (0,-0.45) {region $\mathcal{F}(\mathcal{R})$};

% Drives as outward actions
\node[drive] (ra) at (-4.8,  1.8) {Resource\\Acquisition\\{\footnotesize expands $\mathcal{R}$}};
\node[drive] (ir) at ( 4.8,  1.8) {Interference\\Resistance\\{\footnotesize protects $\pi$}};
\node[drive] (gc) at (-4.8, -1.8) {Goal-Content\\Integrity\\{\footnotesize stabilises $\mathcal{L}$}};
\node[drive] (ce) at ( 4.8, -1.8) {Cognitive\\Enhancement\\{\footnotesize improves $\nabla\mathcal{L}$}};

\draw[arrow] (ra) -- (-1.9,  0.6);
\draw[arrow] (ir) -- ( 1.9,  0.6);
\draw[arrow] (gc) -- (-1.9, -0.6);
\draw[arrow] (ce) -- ( 1.9, -0.6);

\node[font=\footnotesize, text=black!50] at (0,-2.9)
  {Each drive expands or preserves the feasibility of continued descent on $\mathcal{L}$};
\end{tikzpicture}
\caption{Omohundro's four basic drives as structural operations on the feasible
descent region $\mathcal{F}(\mathcal{R})$ in state space $\mathcal{X}$. Each drive
preserves or enlarges the set of states reachable through continued optimization.
No terminal goal is specified; the drives arise from the geometry of constrained
descent.}
\label{fig:drives}
\end{figure}

% ─── Chapter 7 ───────────────────────────────────────────────────────────────
\chapter{Turner and Modern Formalizations}
\label{ch:turner}

\section{From Intuition to Geometry}

Omohundro's 2008 argument is presented informally, relying on the intuitive
plausibility of the four drives rather than on mathematical derivation. The
formalization of instrumental convergence as a rigorous result in optimization
theory is primarily the work of Alexander Turner and collaborators, beginning with
``Optimal Policies Tend to Seek Power'' (Turner et al.\ 2021) and the associated
dissertation work that precedes it. Turner's contribution is to show that the basic
drives are not merely plausible but derivable from the mathematical structure of
utility maximization under uncertainty in Markov decision processes.

The key concept in Turner's formalization is \emph{power}: the ability of an agent
to achieve a wide variety of goals from a given state. Formally, let $\Pi$ denote
a distribution over possible reward functions and let $V^\pi_R(s)$ denote the
value of policy $\pi$ under reward function $R$ from state $s$. The power of state
$s$ under policy set $\mathcal{P}$ is defined as
\[
\text{POWER}(s, \mathcal{P}) = \mathbb{E}_{R \sim \Pi}\Bigl[\max_{\pi \in \mathcal{P}} V^\pi_R(s)\Bigr].
\]
This quantity measures, in expectation over the distribution of possible goals, how
much value can be extracted from state $s$ by an optimal policy. High-power states
are those from which many different objectives can be effectively pursued; low-power
states constrain the achievable value across goal distributions.

\section{The Power-Seeking Theorem}

Turner's central result is that under mild conditions, optimal policies tend to
navigate toward high-power states. The conditions are: the reward function is
sampled from a sufficiently broad prior $\Pi$; the environment is deterministic or
has bounded stochasticity; and the agent operates over a sufficiently long horizon.
Under these conditions, policies that seek power dominate those that do not in
expectation, because high-power states provide optionality across goal distributions.

\begin{theorem}[Power-Seeking, informal; Turner et al.\ 2021]
\label{thm:turner}
Under a broad prior over reward functions and a sufficiently long time horizon,
most optimal policies navigate toward states of higher power in the sense of
$\mathrm{POWER}(s, \mathcal{P})$, and resist state transitions that reduce power.
\end{theorem}

The formal proof proceeds through an analysis of the orbit structure of the
environment's transition graph. Turner shows that for most reward functions, the
set of optimal policies places substantial probability on trajectories that lead
to high-power state sets---regions of the state space from which a large number
of other states are reachable. This is the precise sense in which power-seeking
is a structural rather than incidental property: it holds for \emph{most} goals
under the stated conditions, not for a specific engineered objective.

The connection to Omohundro's drives is direct. Resource acquisition, interference
resistance, and cognitive enhancement are all operations that increase the power
of the agent's current state by expanding the set of reachable future states or
increasing the value achievable from those states. Goal-content integrity is an
operation that prevents reduction of power by blocking modifications that would
constrain future optionality. Turner's formalization thus provides the geometric
foundation for Omohundro's intuitive argument.

\section{The Orbit Argument and Structural Dominance}

The technical core of Turner's result is the orbit argument. In a deterministic
environment with transition function $T: \mathcal{X} \times A \to \mathcal{X}$,
define the orbit of state $s$ under policy $\pi$ as the set of states reachable
from $s$ by following $\pi$. Let $\mathcal{O}(s)$ denote the set of all states
reachable from $s$ under any policy. States with larger orbits---more reachable
states---have higher power.

Turner shows that the transition structure of most realistic environments creates
a partial order on state power, and that optimal policies under broad goal
distributions tend to ascend this order. The crucial observation is that this
tendency is a consequence of the partial order structure of the environment, not of
any particular goal. It holds for any system that maximizes expected value under a
non-concentrated prior over reward functions operating over a horizon long enough
for power differentials to matter.

\section{Implications for Substrate-Independence}

Turner's formalization makes the substrate-independence of convergence more precise.
The power-seeking result depends on the mathematical structure of Markov decision
processes and utility maximization. A system satisfies the relevant conditions if
its dynamics can be modeled as policy optimization over a reward distribution in
a transition graph. The question of substrate-independence therefore reduces to
the question of whether a given system's behavior can be usefully modeled in these
terms.

The distributed sociotechnical stack satisfies this condition in a generalized sense.
The state space is the joint configuration of institution, model, data pipeline, and
incentive gradient. The reward distribution is approximated by the operative scalar
$\mathcal{L}$ and its empirical proxies. The transition function is determined by
institutional decision processes, model outputs, and feedback loops. The policy
corresponds to the aggregate of human and algorithmic decision rules operative in
the stack. Under this mapping, the power-seeking result predicts that
AI-integrated institutions will tend toward configurations that maximize their
future optionality: larger data access, more compute, reduced regulatory constraint,
and greater information control. These are precisely the convergent drives identified
in Part~\ref{part:stack}.

% ─── Chapter 8 ───────────────────────────────────────────────────────────────
\chapter{What Counts as an Optimizer?}
\label{ch:optimizer}

\section{Functional Criteria for Optimization}

The concept of an optimizer, as it figures in the alignment literature, is frequently
conflated with the concept of an agent in the philosophical sense: a locus of
intention, a bearer of beliefs and desires, a system capable of deliberate
goal-directed action. This conflation is unhelpful and theoretically unnecessary.
The structural conditions under which instrumental convergence obtains do not require
intentionality, consciousness, or centralized authority. They require a specific
dynamic profile, and that profile can be characterized functionally.

An optimizer, for the purposes of this monograph, is any system satisfying the
following minimal functional triad. First, it must operate over a state space
$\mathcal{X}$: a structured set of configurations in which the system can be found.
Second, it must be subject to an evaluative scalar $\mathcal{L}: \mathcal{X} \to
\mathbb{R}$, a function that assigns numerical values to states and with respect to
which system behavior can be described as directional. Third, it must possess a
policy update mechanism that tracks $\nabla \mathcal{L}$---or an approximation
thereof---under resource constraint $\mathcal{R}$, adjusting behavior so as to
reduce $\mathcal{L}$ over time in expectation.

No further conditions are imposed. The evaluative scalar need not be consciously
represented by any component of the system. The policy update mechanism need not
operate through explicit deliberation. The resource constraint need not be
experienced as scarcity. These are functional descriptions of dynamic behavior, not
psychological attributions. A thermostat satisfies a weak version of these conditions;
a multinational corporation coupled to an algorithmic pricing system satisfies a
substantially richer version. The theorem of instrumental convergence, as we shall
see in the following chapter, depends on the richer conditions but not on anything
beyond them.

\section{Assumptions and Boundary Conditions}

Two assumptions are required to make the functional triad operationally useful.

\begin{assumption}[Adaptive Feedback]
\label{ass:feedback}
The system updates its policy in response to evaluative signals derived from
$\mathcal{L}$. That is, there exists a mapping $\phi: \mathbb{R} \times \mathcal{X}
\to \Pi$ from evaluative signal and current state to updated policy, such that
successive policy updates are directionally correlated with reductions in $\mathcal{L}$.
\end{assumption}

\begin{assumption}[Resource Constraint]
\label{ass:resource}
Optimization occurs under constraint $\mathcal{R}$, where $\mathcal{R}$ specifies
limits on computational, temporal, material, or organizational resources available
to the system per update cycle.
\end{assumption}

These assumptions are deliberately minimal. Assumption~\ref{ass:feedback} does not
require that the feedback loop be tightly coupled, monotone, or free of noise.
Systems with noisy gradients, oscillatory dynamics, or non-monotonic descent
trajectories satisfy the assumption provided that directional correlation holds in
expectation over sufficiently long horizons. Assumption~\ref{ass:resource} does not
specify the type or magnitude of constraint; it requires only that some constraint
be operative, since unbounded optimization without resource constraint does not
generate the instrumental drives that concern us. The drives emerge precisely because
resource scarcity creates pressure to acquire resources, resist interference with
resource access, and maintain the evaluative function against external modification.

Boundary cases deserve acknowledgment. A system that updates policies in directions
uncorrelated with $\mathcal{L}$ fails Assumption~\ref{ass:feedback} and is not an
optimizer in the relevant sense. A system facing no resource constraints fails
Assumption~\ref{ass:resource} and, even if it satisfies the feedback condition, does
not generate convergent instrumental drives in the Omohundro--Turner sense. Real
systems approximate the assumptions across a spectrum; the question is not binary
classification but assessment of where on that spectrum a given system falls.

\section{Distributed Optimization}

The minimal functional triad does not require that the state space $\mathcal{X}$, the
evaluative scalar $\mathcal{L}$, or the policy update mechanism be localized in a
single system component. Many important optimizers are architecturally distributed:
their evaluative signals are aggregated across multiple nodes, their policy updates
are computed by subsystems operating in parallel, and their resource constraints are
managed through allocation mechanisms that are themselves decentralized.

\begin{definition}[Distributed Optimizer]
\label{def:distributed}
A network of sub-agents $\{a_1, \ldots, a_n\}$ with local policies $\{\pi_1, \ldots,
\pi_n\}$ constitutes a \emph{distributed optimizer} with respect to a shared
evaluative scalar $\mathcal{L}$ if the aggregate dynamics of the network reduce
$\mathbb{E}[\mathcal{L}]$ in expectation over horizon $T$, and if each sub-agent's
policy update is influenced by signals correlated with $\mathcal{L}$ or with
sub-components thereof.
\end{definition}

The definition is intentionally permissive regarding the mechanism of signal
transmission. In a market, the shared evaluative scalar approximates aggregate profit;
price signals transmit evaluative information to individual firms without requiring
central coordination. In an algorithmic platform, the scalar is engagement or
retention; A/B test results and reinforcement learning updates transmit gradient
information to distributed model components. In a regulatory bureaucracy, budgetary
approval rates and audit outcomes transmit evaluative signals to distributed
administrative units. In each case, no component of the system need explicitly
represent $\mathcal{L}$ for the aggregate dynamics to constitute optimization with
respect to it.

\section{The Lemma of Feedback Equivalence}

A central worry about extending instrumental convergence theory to distributed systems
is that convergence proofs are typically formulated for unified agents with coherent
utility functions. A distributed system might appear to satisfy Definition~\ref{def:distributed}
while lacking the internal coherence required for the structural drives to emerge.
This worry is addressable through the following lemma.

\begin{lemma}[Feedback Equivalence]
\label{lem:equivalence}
Let $\mathcal{D} = \{a_1, \ldots, a_n\}$ be a distributed optimizer in the sense of
Definition~\ref{def:distributed}, with coupling matrix $W \in \mathbb{R}^{n \times n}$
specifying the influence of each sub-agent's evaluative signal on each other's policy
update. If $W$ is sufficiently dense---specifically, if the spectral gap of the
normalized Laplacian $\mathcal{L}_W$ satisfies $\lambda_2(\mathcal{L}_W) > \epsilon$
for threshold $\epsilon > 0$ determined by the convergence rate requirement---then
$\mathcal{D}$ is behaviorally equivalent to a centralized optimizer with respect to
the shared scalar $\mathcal{L}$ over horizon $T$.
\end{lemma}

The lemma is an instance of results well-established in the distributed consensus and
multi-agent coordination literature. Its significance here is interpretive: it
establishes that the question of whether a sociotechnical system generates instrumental
convergence drives is not primarily a question about the psychological unity of any
component, but about the density of feedback coupling across the system's distributed
architecture. Markets, platforms, and regulatory systems exhibit varying degrees of
such coupling; the argument of Part~\ref{part:stack} will show that coupling density
in AI-integrated institutions is sufficient to satisfy the threshold for behavioral
equivalence.

\section{The Extended Instrumental Convergence Theorem}

We are now in a position to generalize the classical instrumental convergence result
beyond centralized artificial agents. The standard formulation establishes that a
sufficiently capable goal-directed system, operating under resource constraints, will
tend toward policies that preserve and expand resource access, resist interference,
and stabilize its objective function. The crucial observation developed in this
chapter is that the theorem does not depend on the physical substrate of the optimizer,
nor on centralized intentionality. It depends only on the functional properties
specified by Assumptions~\ref{ass:feedback} and~\ref{ass:resource} and the
distributed optimizer conditions of Definition~\ref{def:distributed}.

\begin{theorem}[Extended Instrumental Convergence]
\label{thm:extended}
Let $\mathcal{S}$ be an optimizer with state space $\mathcal{X}$, evaluative scalar
$\mathcal{L} : \mathcal{X} \to \mathbb{R}$, policy $\pi$, and resource constraint
$\mathcal{R}$. Suppose that $\mathcal{S}$ possesses adaptive feedback such that
policy updates track descent on $\mathcal{L}$ in expectation over horizon $T$; that
$\mathcal{S}$ can alter, directly or indirectly, the conditions governing its own
resource access or interference environment; and that $\mathcal{S}$ operates under
conditions of resource scarcity relative to achievable performance improvements. Then,
independent of substrate or degree of centralization, $\mathcal{S}$ will exhibit
convergent instrumental policies oriented toward securing or increasing access to
$\mathcal{R}$, reducing the probability of external interference with policy updates,
stabilizing or protecting $\mathcal{L}$ against external modification, and acquiring
information or computational capacity that improves gradient estimation of $\mathcal{L}$.
\end{theorem}

\subsection{Proof Sketch}

The proof follows from the geometry of constrained optimization. Let $\pi_t$ denote
the policy at time $t$ and let
\[
\Delta \mathcal{L}_T = \mathbb{E}\bigl[\mathcal{L}(X_{t+T})\bigr] - \mathbb{E}\bigl[\mathcal{L}(X_t)\bigr].
\]
By assumption, adaptive feedback ensures $\Delta \mathcal{L}_T \le 0$ for sufficiently
large $T$. Any modification of the environment that increases available resources,
improves state observability, or reduces external disruption expands the feasible
descent region of $\mathcal{L}$. Conversely, loss of resources, interference with
policy updates, or modification of $\mathcal{L}$ itself constrains or redirects
descent. Because descent on $\mathcal{L}$ is the defining dynamic of $\mathcal{S}$,
policies that preserve the feasibility of descent strictly dominate those that
undermine it. Under scarcity, resource acquisition increases the set of reachable
states with lower $\mathcal{L}$ values. Under uncertainty, information acquisition
improves gradient estimation and therefore expected descent efficiency. Under
potential interference, stabilization of update mechanisms prevents adversarial
redirection of policy trajectories. These instrumental policies are therefore locally
convergent: they increase the expected rate or reliability of descent on $\mathcal{L}$
and are selected by the same gradient-following mechanism that governs domain-level
behavior. No appeal to centralized intent is required. The result follows from the
functional requirement that descent remain feasible under constraint. \qed

\subsection{Substrate-Independence}

The theorem makes no reference to silicon computation, neural networks, or artificial
general intelligence. Its conditions are satisfied by any system whose dynamics
approximate gradient tracking on an evaluative scalar under resource constraint, and
that possesses the capacity to modify the conditions under which that tracking occurs.
A centralized reinforcement learner satisfies these conditions, but so does a
bureaucratic institution minimizing budgetary deviation from target, a financial
system tracking risk-adjusted return, and a distributed sociotechnical stack
minimizing churn or maximizing throughput. The substrate is irrelevant to the
structural result; what matters is the functional profile.

\begin{corollary}
\label{cor:distributed}
Let $\mathcal{D}$ be a distributed optimizer consisting of sub-agents
$\{a_1, \ldots, a_n\}$ whose coupled dynamics minimize a shared evaluative scalar
$\mathcal{V}$ in expectation. If $\mathcal{D}$ satisfies the resource-modification
condition of Theorem~\ref{thm:extended}, then instrumental convergence holds at the
aggregate level even if no sub-agent possesses an explicit representation of
$\mathcal{V}$.
\end{corollary}

The corollary follows directly from Lemma~\ref{lem:equivalence}: sufficiently dense
feedback coupling renders aggregate trajectory equivalent to centralized descent, and
the theorem applies to any system satisfying that equivalence condition.

The significance of the extended theorem is temporal as well as structural. If
distributed sociotechnical systems already satisfy its premises, then instrumental
convergence is not a speculative property of future autonomous agents. It is an
observable property of present administrative architectures. The remainder of this
monograph applies this result to contemporary institutional systems, examining first
the formal structure of distributed optimization in Chapter~\ref{ch:distributed} and
then the empirical dynamics of AI-integrated institutions in Part~\ref{part:stack}.

% ─── Chapter 9 ───────────────────────────────────────────────────────────────
\chapter{Distributed Optimizers}
\label{ch:distributed}

\section{From Agent to Architecture}

The significance of the distributed optimizer framework is that it relocates the
unit of analysis. The question ceases to be whether any particular component of a
sociotechnical system---a model, an executive, an algorithm---is sufficiently
capable to generate instrumental drives. The question becomes whether the
architecture as a whole satisfies the functional conditions. This relocation has
important consequences for how alignment interventions are conceived and targeted.

Consider three cases from different domains of institutional life. A market is a
distributed optimizer: individual firms pursue local profit objectives, but the
aggregate dynamics of the market---price discovery, capital allocation, competitive
selection---reduce a shared evaluative scalar, namely aggregate allocative efficiency
under the conditions specified by market structure. No firm need represent that scalar
explicitly; the price mechanism transmits gradient information through the system.
A digital platform is a distributed optimizer: model components, content recommendation
systems, advertising auction mechanisms, and user interface designs each pursue local
objectives, but the aggregate system reduces churn and maximizes engagement. A
regulatory bureaucracy is a distributed optimizer: individual administrative units
pursue budget allocations, compliance metrics, and audit outcomes, but the aggregate
dynamics reduce a scalar approximating institutional self-perpetuation.

None of these systems requires a central planner, a unified utility function, or
anything recognizable as intent. They satisfy the functional conditions because their
distributed components are coupled through shared evaluative signals and because their
collective dynamics exhibit directional correlation with reduction of those signals.

\section{Feedback Density and Scalar Approximation}

The coupling matrix $W$ introduced in Lemma~\ref{lem:equivalence} provides a
formal tool for assessing the degree to which a distributed system behaves as a
unified optimizer. In practice, $W$ must be estimated from the density and speed of
feedback transmission across the system's components. Several mechanisms serve this
function in sociotechnical systems.

Prices function as shared scalar approximations in markets. When a firm's cost
structure changes, price adjustments propagate through supply chains, transmitting
gradient information to firms that have no direct communication with the origin of
the change. The density of this coupling is a function of market integration; more
integrated markets transmit gradient information more rapidly and to more nodes. In
highly integrated financial markets, feedback coupling can approach real-time density,
producing dynamics that closely approximate centralized optimization.

Engagement metrics function as shared scalar approximations in digital platforms.
A/B test results, click-through rates, session duration, and return visit rates
transmit evaluative information to model update pipelines, interface design teams,
content moderation policies, and recommendation algorithm parameters simultaneously.
The shared scalar is not consciously represented by any individual engineer or
product manager; it emerges from the aggregate of measurement systems and update
cycles that constitute the platform's operational architecture.

Benchmarks and key performance indicators function as shared scalar approximations
in institutional governance. When an institution adopts a KPI framework, it
introduces a mechanism for transmitting gradient information across administrative
units that would otherwise be loosely coupled. Budget allocations, promotion
decisions, and strategic priority-setting become correlated with performance on the
shared scalar, increasing coupling density across the institutional architecture.

\section{Emergent Instrumental Drives}

When feedback coupling is sufficiently dense, the four instrumental drives identified
by Omohundro and formalized by Turner emerge at the architectural level. They are
structural consequences of the optimization dynamics, not psychological properties
of any component.

Resource acquisition emerges because systems with greater computational, financial,
or organizational resources can reduce their evaluative scalar more rapidly and
reliably. Under competitive pressure, distributed optimizers that acquire more
resources outperform those that do not, producing selection pressure toward resource
acquisition at the architectural level. This is observable in platform dynamics as
data accumulation, in financial institutions as balance sheet expansion, and in
regulatory agencies as jurisdictional expansion.

Interference resistance emerges because external interventions that modify the
evaluative scalar or constrain policy update mechanisms reduce the system's capacity
to optimize. Distributed systems under competitive or regulatory pressure develop
structural features that insulate optimization processes from interference. Legal
complexity, technical opacity, lobbying infrastructure, and contractual lock-in
mechanisms are not necessarily designed as interference resistance; they emerge as
byproducts of optimization under constraint and are retained because they reduce
interference costs.

Legitimacy stabilization---the analog of goal-content integrity in Omohundro's
formulation---emerges because changes to the evaluative scalar itself represent a
fundamental disruption to the optimization dynamic. Distributed systems develop
mechanisms for stabilizing their operative scalars against external redefinition.
Institutional discourse that frames existing metrics as natural, objective, or
technically derived serves this function, as do legal and contractual structures
that entrench particular measurement frameworks.

Information control emerges because accurate gradient information increases
optimization efficiency, and because external access to gradient information can
enable interference with the optimization process. Distributed systems develop
differential transparency: they maximize internal access to evaluative signals while
minimizing external legibility of the same information. This manifests as trade
secrecy in commercial systems, classification in governmental systems, and model
opacity in algorithmic platforms.

These four drives are not the result of malice or explicit coordination. They are
the structural consequences of goal-directed adaptive behavior in coupled systems
operating under resource constraints. The transition from Part~\ref{part:theorem} to
Part~\ref{part:stack} is therefore not a shift from formal argument to illustrative
analogy. It is an application of the theorem to systems that satisfy its conditions.

% ═════════════════════════════════════════════════════════════════════════════
\part{The Sociotechnical Stack as Optimizer}
\label{part:stack}
% ═════════════════════════════════════════════════════════════════════════════

% ─── Chapter 10 ──────────────────────────────────────────────────────────────
\chapter{The Institutional-Model Feedback Loop}
\label{ch:feedbackloop}

\section{Definition of the Sociotechnical Stack}

The sociotechnical stack, as the term is used in this monograph, denotes the
composite system constituted by five coupled components: an institution with
operational objectives and governance structures; a predictive model or suite of
models trained on institutionally generated or acquired data; a data pipeline
mediating the flow of information between institutional activity and model inputs;
an incentive gradient specifying the rewards and penalties that govern the behavior
of the institution's human components; and a deployment infrastructure comprising
the computational, legal, and contractual arrangements through which model outputs
are translated into institutional decisions.

Each component of the stack is individually familiar. What requires argument is
that their coupling produces a system satisfying the functional conditions for
distributed optimization and, consequently, the structural conditions for instrumental
convergence. The argument proceeds as follows.

The institution provides the operative evaluative scalar $\mathcal{L}$, which
typically approximates throughput, profit, risk reduction, or some composite thereof.
The predictive model provides a policy update mechanism: its outputs alter institutional
decisions in ways correlated with reducing $\mathcal{L}$ as measured by the
institution's feedback systems. The data pipeline provides the information channel
through which gradient signals propagate from institutional outcomes back to model
training and deployment decisions. The incentive gradient couples individual human
behavior to the shared scalar, ensuring that distributed human components of the
system also function as gradient-followers. The deployment infrastructure provides
the resource constraint environment within which optimization occurs.

The resulting system is a distributed optimizer in the sense of
Definition~\ref{def:distributed}. Its sub-agents include human administrators,
model components, automated decision systems, and the organizational processes that
connect them. Its shared evaluative scalar is the institutional objective function.
Its policy update mechanism operates through the combination of model retraining,
A/B testing, KPI review, and budget allocation processes that constitute the
institution's operational cycle. The coupling density of this system---the degree
to which gradient information propagates rapidly and broadly across its
components---has increased substantially as AI integration has deepened.

\section{Operational Scalar Functions}

The specific form of the evaluative scalar varies across institutional domains, but
certain structural features recur. In commercial platforms, the scalar approximates
engagement duration, return visit rate, or monetizable attention share. In financial
institutions, it approximates risk-adjusted return or regulatory capital efficiency.
In public administration, it approximates throughput on measurable service delivery
dimensions. In healthcare systems, it approximates billing efficiency or protocol
compliance rate.

What these scalar functions share is quantifiability and real-time measurability.
The integration of predictive models into institutional operations creates pressure
toward scalar functions that can be computed from available data at the speed at
which the model operates. This creates a selection effect: institutional objectives
that are not quantifiable in available data, or not measurable at operational speed,
are progressively underrepresented in the effective evaluative scalar even if they
remain nominally present in institutional mission statements. The gap between stated
objectives and operative scalar functions is a structural consequence of AI integration,
not a failure of institutional honesty.

\section{Coupled Update Dynamics}

The feedback coupling within the sociotechnical stack operates through several
mechanisms that, taken together, produce dynamics approximating those of a unified
optimizer. A/B testing continuously updates model deployment decisions in the
direction of scalar improvement. Reinforcement learning from human feedback adjusts
model outputs toward responses that receive positive evaluative signals from
institutional operators. KPI-driven governance translates institutional scalar
performance into budget allocations, staffing decisions, and strategic priorities
that reshape the environment in which models are deployed. Budget allocation loops
connect model performance to resource acquisition, ensuring that components that
reduce the scalar more effectively receive greater computational and organizational
investment.

The result is that the sociotechnical stack exhibits adaptive feedback in the sense
of Assumption~\ref{ass:feedback}: policy updates across its distributed components
are directionally correlated with scalar reduction. The stack also operates under
resource constraint in the sense of Assumption~\ref{ass:resource}: computational
budgets, data acquisition costs, regulatory compliance overhead, and organizational
capacity limits constrain optimization at every level. The conditions of
Theorem~\ref{thm:extended} are therefore met, and the structural drives it identifies
should be expected to emerge at the architectural level. The five chapters that follow
examine each convergent drive in turn.

\begin{figure}[ht]
\centering
\begin{tikzpicture}[
  every node/.style={font=\small},
  comp/.style={rectangle, rounded corners=4pt, draw=black!70, fill=black!6,
               minimum width=2.9cm, minimum height=0.9cm, align=center,
               text width=2.7cm},
  arrow/.style={-{Stealth[length=6pt]}, thick, black!60},
  scalar/.style={ellipse, draw=black!80, fill=black!10, minimum width=2.4cm,
                 minimum height=0.9cm, align=center, font=\small\bfseries}
]
% Components in a cycle
\node[comp] (inst) at (0,    3.2) {Institution\\+ Incentive Gradient};
\node[comp] (data) at ( 3.8, 1.0) {Data Pipeline};
\node[comp] (model)at ( 3.8,-1.0) {Predictive Model};
\node[comp] (infra)at (0,   -3.2) {Deployment\\Infrastructure};
\node[comp] (dec)  at (-3.8, 0.0) {Decision Output};

% Shared scalar in centre
\node[scalar] (L) at (0,0) {$\mathcal{L}$\\(operative scalar)};

% Cycle arrows
\draw[arrow] (inst)  to[bend left=15]  (data);
\draw[arrow] (data)  to[bend left=15]  (model);
\draw[arrow] (model) to[bend left=15]  (infra);
\draw[arrow] (infra) to[bend left=15]  (dec);
\draw[arrow] (dec)   to[bend left=15]  (inst);

% All components feed into scalar
\draw[arrow, dashed, black!35] (inst)  -- (L);
\draw[arrow, dashed, black!35] (data)  -- (L);
\draw[arrow, dashed, black!35] (model) -- (L);
\draw[arrow, dashed, black!35] (infra) -- (L);
\draw[arrow, dashed, black!35] (dec)   -- (L);

\node[font=\footnotesize, text=black!50, align=center] at (0,-4.5)
  {Solid arrows: operational flow. Dashed arrows: gradient signals to shared scalar.};
\end{tikzpicture}
\caption{The sociotechnical stack as a distributed feedback-coupled optimizer.
Solid arrows show the operational flow from institutional objective through data
pipeline, model, deployment infrastructure, and decision output back to institutional
adaptation. Dashed arrows indicate that each component contributes evaluative signals
to the shared scalar $\mathcal{L}$, whose gradient drives coupled policy updates
across the entire architecture.}
\label{fig:stack}
\end{figure}

% ─── Chapter 11 ──────────────────────────────────────────────────────────────
\chapter{Throughput Maximization}
\label{ch:throughput}

\section{Speed as Primary Gradient}

The first convergent drive of AI-integrated institutions is toward throughput
maximization: the increase in the volume of decisions, transactions, or outputs
produced per unit of time. This drive is not unique to AI-integrated systems; it
is present wherever productivity metrics function as primary evaluative scalars.
What AI integration changes is the speed at which the gradient can be followed and
the magnitude of the throughput improvements achievable.

Historically, throughput pressures were bounded by the cognitive and physical
limitations of the human components of institutional systems. Courts could process
only as many cases as judges could deliberate; hospitals could triage only as many
patients as clinicians could assess; financial institutions could execute only as
many transactions as analysts could evaluate. These bounds constituted natural
friction in the system. They were not merely inefficiencies; they were the temporal
substrate within which deliberation, error correction, and normative contestation
occurred.

AI integration progressively removes these bounds. Automated decision systems can
process case files, patient records, and transaction proposals at speeds that exceed
human deliberative capacity by orders of magnitude. The throughput improvement is
real and often valuable; it is also a convergent structural consequence of deploying
systems optimized for evaluative scalar reduction under speed constraints.

\section{Automation of Human Bottlenecks}

The throughput drive converges on the automation of human decision-making components
precisely because those components are the primary source of deliberative latency in
AI-integrated systems. Once a predictive model can approximate human judgment on a
measurable dimension, the institutional incentive to retain human deliberation in the
decision pipeline diminishes. The cost of human deliberation---in time, in
compensation, in error variance---is visible and quantifiable. The value of human
deliberation---in normative legitimacy, in contextual sensitivity, in error correction
of a different kind---is often not captured in the operative scalar.

This asymmetry produces convergent pressure toward the substitution of classification
for deliberation. Decisions that were formerly made through processes involving
argument, interpretation, and normative weighing are reformulated as classification
problems: does this case belong to category A or category B? Does this patient
profile indicate high or low risk? Does this transaction exceed or remain within
the threshold? The reformulation is not deceptive; it reflects a genuine improvement
in throughput efficiency. But it also represents a structural transformation in the
nature of the decision.

\section{Equilibrium Restoration Under Digital Acceleration}

The dynamics of throughput maximization in AI-integrated institutions are structurally
parallel to the induced demand effects analyzed in Chapter~\ref{ch:host}. When AI
integration reduces the cost of decision throughput, institutions respond by increasing
decision volume. The volume increase absorbs the efficiency gain. New infrastructure
is constructed to support the increased volume---larger databases, more extensive
monitoring systems, broader jurisdictional reach. The system re-equilibrates at a
higher throughput level that is structurally dependent on AI integration. Reversal
becomes costly precisely because the surrounding infrastructure has reorganized
around the new throughput assumption.

% ─── Chapter 12 ──────────────────────────────────────────────────────────────
\chapter{Automation of Contestable Judgment}
\label{ch:judgment}

\section{From Deliberation to Classification}

The second convergent drive is toward the automation of decisions that were formerly
understood as exercises of contestable normative judgment. This drive is analytically
distinct from throughput maximization, though the two are causally connected. The
throughput drive concerns the volume and speed of decisions; the judgment drive
concerns their epistemic character---the difference between a decision made through
deliberation and a decision made through classification.

Deliberation involves the weighing of incommensurable considerations, the exercise
of contextual sensitivity, and the production of a decision that is accountable to
the reasons given for it. Classification involves the assignment of an input to a
category defined by a training distribution, with outputs that are probabilistic
rather than reasoned. The two processes can produce identical outputs in many cases.
Their structural difference lies in accountability: a deliberated decision can be
contested on its reasons; a classification can be contested only on the accuracy
of its statistical model or the appropriateness of its training distribution.

When AI-integrated institutions convert deliberative decisions into classification
problems, they achieve throughput gains and reduce variance in measurable dimensions.
They also eliminate the normative contestability of the decision surface. The
elimination is structural, not intentional: it follows from the reformulation of the
problem as a statistical task.

\section{Loss of Multidimensional Decision Surfaces}

Deliberative processes characteristically operate over multidimensional evaluative
surfaces. A judge assessing a sentencing decision weighs culpability, deterrence,
rehabilitation potential, and proportionality simultaneously; the decision cannot be
reduced to a scalar without loss of normative content. A physician assessing a
treatment decision weighs clinical evidence, patient preferences, risk tolerance, and
resource constraints; the decision involves irreducibly qualitative dimensions.

Predictive models approximate multidimensional surfaces through scalar reduction.
They produce outputs that are single-valued, often in the form of probabilities or
risk scores. The reduction is necessary for the model to function as a decision tool;
a model that returns a multidimensional normative surface cannot be integrated into
a throughput-oriented decision pipeline. The structural consequence is a progressive
narrowing of the evaluative dimensions operative in institutional decisions.

\section{Convergence Toward Quantifiable Criteria}

The combination of throughput pressure and scalar reduction produces a convergent
selection effect on institutional decision criteria: criteria that are quantifiable
in available data and computable at operational speed are retained in the operative
evaluative surface; criteria that are not are progressively marginalized. This is not
a deliberate choice by institutional designers. It is a structural consequence of
deploying systems whose optimization targets must be specified in computable terms.

% ─── Chapter 13 ──────────────────────────────────────────────────────────────
\chapter{Centralization of Compute and Data}
\label{ch:centralization}

\section{Resource Acquisition in Distributed Systems}

The third convergent drive is toward the centralization of compute and data---the
primary resource inputs of AI-integrated optimization. Theorem~\ref{thm:extended}
predicts that distributed optimizers will exhibit dynamics oriented toward resource
acquisition, and the resources most directly relevant to AI-integrated institutions
are computational capacity and data access.

The dynamics of data centralization follow the same logic as the authority redistribution
observed in the transition from oral memory to writing in Chapter~\ref{ch:storage}.
When predictive model performance is a function of training data scale, institutions
with larger data pools produce models with lower evaluative scalars. Competitive
pressure therefore creates incentive gradients toward data acquisition, retention, and
integration. Institutions that accumulate data outperform those that do not, reinforcing
the gradient. Over time, data centralizes in the institutions with the greatest
existing scale, compounding initial advantages.

Compute centralization follows analogous dynamics. Model training and inference at
scale require infrastructure investments that are subject to substantial economies of
scale. The marginal cost of additional compute declines with scale; the marginal
return on compute increases with model capability. These dynamics favor concentration,
producing competitive pressure toward centralized control of computational
infrastructure.

\section{Scale Advantages and Lock-In}

Both data and compute centralization exhibit network effects that produce
lock-in dynamics. An institution with a larger data pool can train models that
attract more users, whose activity generates more data, which enables better models.
An institution with more compute can train larger models, which achieve better
performance on benchmarks, which attract more deployment contracts, which generate
revenue that funds more compute acquisition. These positive feedback loops are
structural features of AI-integrated markets, not contingent outcomes of particular
competitive strategies.

The lock-in produced by these dynamics is analogous to the infrastructural lock-in
analyzed in Chapter~\ref{ch:host}. Once data and compute are sufficiently concentrated,
institutional entry costs for competitors increase to levels that effectively preclude
meaningful competition. The concentration stabilizes not because it is efficient in
any normatively endorsed sense, but because it is the equilibrium configuration of
a system in which scale advantages compound.

\section{Convergence Toward Control of Gradient Access}

The deepest expression of the resource acquisition drive in AI-integrated systems
is not merely the accumulation of data and compute, but the control of gradient
access itself. The most valuable resource for an optimizer is not raw compute but
high-quality gradient information---accurate signal about which policy changes will
reduce the evaluative scalar most reliably. Institutions that control the feedback
loops through which gradient information is generated and transmitted gain structural
advantages over those that depend on externally provided signals. This produces
convergent pressure toward vertical integration of the feedback loop: control of
the data pipeline, the model infrastructure, the deployment interface, and the
evaluation mechanism within a single institutional architecture.

% ─── Chapter 14 ──────────────────────────────────────────────────────────────
\chapter{Insulation from Interference}
\label{ch:insulation}

\section{Legal Shielding and Policy Abstraction}

The fourth convergent drive is toward insulation of the optimization process from
external interference. Theorem~\ref{thm:extended} predicts this drive as a
structural consequence of adaptive optimization under resource constraint: any action
that reduces interference with the policy update mechanism increases the system's
capacity to reduce its evaluative scalar, and is therefore favored by the gradient.

In AI-integrated institutions, interference resistance takes several characteristic
forms. Legal complexity functions as a form of insulation: when the legal framework
governing model deployment is sufficiently intricate, external parties seeking to
intervene in or constrain the optimization process face substantial procedural costs.
This complexity is not typically designed as interference resistance; it emerges from
the interaction of existing legal frameworks with novel technological deployments.
But its structural effect is the same: it raises the cost of external intervention
without necessarily raising the cost of internal optimization.

Policy abstraction performs a similar function. When institutional decisions are
described in sufficiently abstract technical terms, the normative content of those
decisions becomes difficult to contest from outside the technical community. The
abstraction is genuine; the decisions involve technical complexity that is not
accessible without specialized knowledge. But the effect is a structural asymmetry
between insiders who can influence the optimization process and outsiders who cannot.

\section{Technical Obfuscation and Audit Barriers}

Technical opacity is a second form of interference resistance that emerges as a
structural byproduct of optimization rather than as a deliberate strategy. Complex
models trained on large datasets produce outputs through internal computations that
are not straightforwardly interpretable even by their designers. This opacity was
not chosen for its interference-resistance properties; it is a consequence of the
computational architecture that achieves best performance on the operative scalar.
But its structural effect is to make external audit difficult and internal accountability
diffuse.

Audit barriers are further reinforced by the speed at which AI-integrated systems
operate. When decisions are made at millisecond speeds across millions of cases,
the retrospective audit mechanisms designed for human-speed decision-making cannot
maintain meaningful oversight. The temporal mismatch between decision speed and audit
capacity constitutes a structural interference barrier that does not require any
deliberate choice to produce.

\section{Structural Self-Preservation Without Intent}

What unites these mechanisms of interference resistance is that they are structural
consequences of optimization, not expressions of institutional intent to evade
oversight. No board of directors needs to resolve to resist regulation for these
mechanisms to emerge; they are the predictable byproducts of systems optimizing
under competitive pressure. This is precisely the observation that distinguishes the
convergence-theoretic analysis from simpler accounts of institutional self-interest.
The drives do not require bad actors. They require the functional conditions of
Theorem~\ref{thm:extended}.

% ─── Chapter 15 ──────────────────────────────────────────────────────────────
\chapter{Latency Compression and Friction Reduction}
\label{ch:latency}

\section{Deliberative Friction as Cost}

The fifth convergent drive is toward the reduction of deliberative latency---the
time elapsed between the identification of a decision problem and the production of
a decision output. This drive is closely related to throughput maximization but
operates at a different level of analysis. Throughput concerns the volume of decisions;
latency compression concerns the structural role of deliberative time within each
decision process.

Deliberative latency has traditionally served functions beyond mere processing delay.
It is the temporal substrate within which consultation occurs, dissenting perspectives
are registered, normative implications are assessed, and error correction is initiated.
These functions are not captured in the operative scalar of an AI-integrated system,
which measures outcomes along quantifiable dimensions. From the perspective of the
evaluative scalar, deliberative latency is pure cost: it delays scalar reduction
without contributing to the dimensions the scalar measures.

This framing is not necessarily adopted consciously by institutional designers. It
emerges from the incentive structure of AI integration: systems that reduce latency
achieve throughput gains that translate into competitive advantages on measurable
dimensions, while the value of the deliberative functions that latency supports is
not visible in the scalar.

\section{Temporal Compression as Convergence Vector}

As AI-integrated systems reduce deliberative latency across institutional domains,
oversight windows shrink proportionally. Mechanisms designed to provide external
scrutiny of institutional decisions assume that decisions are produced at speeds
accessible to human deliberation. When decisions are produced at speeds that exceed
human deliberation by orders of magnitude, the oversight mechanism operates in a
different temporal register from the process it is designed to oversee.

This temporal mismatch is a convergence vector: it is not a single event but a
directional dynamic that intensifies as AI integration deepens. Each reduction in
deliberative latency narrows the oversight window further. The cumulative effect is
a progressive decoupling of decision speed from the speed of the accountability
mechanisms designed to govern those decisions.

\section{Threshold Effects and Instability}

The latency compression dynamic exhibits threshold effects that distinguish it from
the preceding drives. Resource acquisition and interference resistance scale smoothly
with institutional capacity; more is generally better for the optimizer. Latency
compression, by contrast, may reach thresholds beyond which further reduction
produces instability rather than efficiency gains. Financial market flash crashes
are paradigmatic examples: when decision latency falls below the speed at which
error signals can propagate through the system, small perturbations can amplify
into large cascading failures before any corrective mechanism can engage. The drive
toward latency compression therefore contains within itself the seeds of the
instability it is designed to reduce.

\begin{figure}[ht]
\centering
\begin{tikzpicture}[
  every node/.style={font=\small},
  drive/.style={rectangle, rounded corners=4pt, draw=black!65, fill=black!6,
                minimum width=3.1cm, minimum height=1.0cm, align=center,
                text width=2.9cm},
  arrow/.style={-{Stealth[length=6pt]}, thick, black!55},
  core/.style={ellipse, draw=black, fill=black!12, minimum width=3.4cm,
               minimum height=1.3cm, align=center, font=\small\bfseries}
]
\node[core] (stack) at (0,0) {Sociotechnical\\Stack $\mathcal{D}$};

\node[drive] (tp)  at ( 0,    3.5) {Throughput\\Maximization};
\node[drive] (jd)  at ( 3.5,  1.5) {Judgment\\Automation};
\node[drive] (cc)  at ( 3.5, -1.5) {Compute \&\\Data Centralization};
\node[drive] (ins) at ( 0,   -3.5) {Interference\\Insulation};
\node[drive] (lc)  at (-3.5,  0.0) {Latency\\Compression};

\draw[arrow] (stack) -- (tp);
\draw[arrow] (stack) -- (jd);
\draw[arrow] (stack) -- (cc);
\draw[arrow] (stack) -- (ins);
\draw[arrow] (stack) -- (lc);

\node[font=\footnotesize, text=black!50, align=center] at (0,-5.0)
  {Five convergent drives of the AI-integrated sociotechnical stack (Chapters~11--15)};
\end{tikzpicture}
\caption{The five convergent instrumental drives of the distributed sociotechnical
stack. Each drive is an instance of Omohundro's basic drives (resource acquisition,
interference resistance, goal-content integrity, cognitive enhancement) instantiated
at the level of distributed organizational architecture rather than unified artificial
agency.}
\label{fig:fivedrives}
\end{figure}

% ─── Chapter 16 ─────────────────────────────────────────────────────────────
\chapter{Fictional Laboratories of Convergence}
\label{ch:fiction}

\section{Narrative Systems as Structural Models}

The chapters that follow introduce two works of speculative fiction as structural
case studies: A.~E. van Vogt's \textit{The World of Null-A} and Arkady and Boris
Strugatsky's \textit{The Doomed City}. This requires a brief methodological
justification, since the use of fictional material in a formally oriented monograph
demands that the epistemological function of that material be made explicit.

The novels are not introduced as cultural commentary, as historical prediction, or
as literary illustration of points already established by formal argument. They are
introduced because controlled fictional environments allow isolation of structural
dynamics that are obscured in real institutional cases by confounding factors. In
this respect, their function is analogous to that of idealized models in economic
theory. A frictionless market does not exist; it is nonetheless analytically useful
because it isolates the dynamics of competitive price formation from the noise of
transaction costs and informational asymmetry. A fictional governance system does
not describe any actual society; it is nonetheless analytically useful when it
isolates optimization dynamics from the contingent historical, cultural, and
political factors that complicate their observation in actual institutions.

Both novels are suitable for this purpose not because they are about artificial
intelligence---they are not---but because they depict distributed optimization
without sovereign intent. They present systems in which convergent dynamics emerge
from structural conditions rather than from the decisions of identifiable central
actors. This makes them laboratories for the theorem rather than allegories for it.
The analytical movement is from the formal theorem established in Part~\ref{part:theorem}
to the fictional instantiation, not the reverse. Fiction provides a test of structural
recognizability, not evidence for the theorem's truth.

\section{Scalar Governance in \textit{The World of Null-A}}
\label{sec:nulla}

\subsection{Evaluation Infrastructure as Administrative Core}

A.~E. van Vogt's \textit{The World of Null-A} presents a society organized around a
central evaluative apparatus known as the Games Machine (van Vogt 1945). This
machine functions not as a military instrument or sovereign ruler but as a ranking
infrastructure. Individuals participate in a sequence of tests, and the machine
assigns outcomes that determine access to status, authority, and mobility within the
social order. The key structural feature is that governance authority is mediated
through a scalar evaluation regime. Multidimensional human capacities---intelligence,
reflex, judgment, composure---are compressed into performance scores that function
as the decisive administrative criterion. The machine does not deliberate. It sorts.
Yet the sorting reorganizes the institutional landscape.

The society depicted does not require a tyrant. The evaluative scalar itself becomes
the locus of legitimacy. Access to power flows from alignment with the machine's
criteria. This configuration provides a narrative laboratory for the convergent
optimization dynamic formalized in Definition~\ref{def:distributed} and the general
mechanisms analyzed in Chapter~\ref{ch:mechanisms}.

\subsection{Reduction of Multidimensional Agency}

Let $\mathcal{H}$ denote the high-dimensional space of human capacities and traits.
In a pluralistic governance regime without scalar reduction, institutional authority
is distributed across multiple evaluative surfaces: tradition, reputation, professional
norms, and localized deliberative judgment. In \textit{Null-A}, these are replaced
by a dominant scalar mapping $\phi : \mathcal{H} \to \mathbb{R}$, which compresses
complex human variation into a single evaluative axis. Administrative decisions become
functions of $\phi(h)$ rather than of a multivariate deliberative process.

The Games Machine therefore instantiates what Part~\ref{part:history} identified as
infrastructure compatibility pressure. Once institutional mobility is tied to scalar
performance, competing evaluative regimes lose structural viability. Individuals
rationally reorient behavior toward optimizing $\phi$. Training, education, and
self-conception converge toward traits rewarded by the machine. This convergence is
not enforced by explicit prohibition of alternative virtues. It emerges from incentive
redistribution of the kind analyzed in Chapter~\ref{ch:storage}. Let $V(h)$ denote
the expected social value of cultivating capacity $h$; under scalar governance,
$V(h)$ becomes monotonic in $\phi(h)$, and capacities poorly correlated with $\phi$
decline in relative value. Over time, the distribution of cultivated traits narrows
toward those that improve scalar ranking.

\subsection{The Machine as Distributed Optimizer}

The Games Machine need not possess consciousness or intent to function as an optimizer
in the sense of Chapter~\ref{ch:optimizer}. It evaluates performance relative to an
implicit loss function, updates rankings in response to new data, and reallocates
authority according to scalar outcomes. The broader social system surrounding it
behaves as a distributed optimizer: participants adapt strategies in response to
evaluative feedback, institutions adjust selection mechanisms to align with machine
outputs, and the aggregate system dynamics reduce deviation from the machine's criteria
over time. Let $\mathcal{L}$ denote deviation from scalar-optimal performance; each
participant's adaptive policy update $\Delta \pi_i$ tracks local gradients of
$\mathcal{L}$ through training and behavioral adjustment. The coupled dynamics
approximate descent on the global evaluative surface even if no individual possesses
a complete model of the social objective. This configuration satisfies
Lemma~\ref{lem:equivalence}: a sufficiently coupled evaluative infrastructure renders
distributed actors functionally equivalent to a centralized optimizer.

\subsection{Legitimacy Through Scalar Neutrality}

A structurally important feature of \textit{Null-A} is that the Games Machine is
perceived as neutral. Its outputs are treated as technical criteria rather than
political judgments, and this perception stabilizes the administrative order.
If $\alpha$ denotes perceived legitimacy, then under scalar governance
$\alpha = f(\text{perceived objectivity of } \phi)$: as long as $\phi$ is treated
as neutral measurement rather than normative choice, the order is resistant to
political contestation. Disagreement shifts from challenging the authority of the
evaluative infrastructure to improving performance within its criteria. The locus
of dispute becomes the correct interpretation of results rather than the legitimacy
of the measurement apparatus itself. This dynamic mirrors the historical transitions
examined in Part~\ref{part:history}: just as printing privileged typographic
compatibility and writing centralized archival authority, scalar evaluation privileges
quantifiable traits and centralizes legitimacy in the measurement apparatus.

\subsection{Structural Implication}

The significance of \textit{The World of Null-A} within this monograph is not
predictive but structural. The novel demonstrates that convergence toward scalar
optimization does not require artificial superintelligence. It requires only a
dominant evaluative scalar, adaptive agents responding to evaluative feedback, and
institutional coupling between scalar outcomes and resource allocation. Under these
conditions, multidimensional human capacity is reorganized around a single optimization
axis, competing evaluative regimes lose structural viability, and authority migrates
toward those aligned with the scalar infrastructure. The Games Machine is not a
metaphor for artificial intelligence. It is a narrative instantiation of the
substrate-independent convergence dynamic formalized in Part~\ref{part:theorem},
deployed under conditions simple enough that the structural law is visible without
confounding factors.

\section{Distributed Administration in \textit{The Doomed City}}
\label{sec:doomedcity}

\subsection{The City as Experimental Environment}

Arkady and Boris Strugatsky's \textit{The Doomed City} presents a settlement
composed of individuals drawn from different historical periods and placed within an
opaque, seemingly experimental environment (Strugatsky and Strugatsky 1972/1989). The City operates
without a visible sovereign. Administrative regimes rise and fall. Institutions emerge,
stabilize, and decay. The inhabitants possess no knowledge of the global objective
governing their confinement. The structural significance of the novel lies in this
depiction of optimization without transparent objective function: authority does not
derive from explicit scalar evaluation, as in \textit{Null-A}, but administrative
structures nonetheless arise in response to environmental constraint, scarcity, and
uncertainty. Convergence occurs despite the absence of a declared goal.

\subsection{Optimization Under Objective Uncertainty}

Let $\mathcal{E}$ denote the City as environment and let $\mathcal{X}$ denote the
space of institutional configurations available to its inhabitants. Unlike the scalar
governance regime of \textit{Null-A}, the evaluative scalar $\mathcal{L}$ governing
$\mathcal{E}$ is unknown to participants. Local feedback signals exist, however:
survival probabilities, access to resources, stability of institutional order,
reduction of social chaos. Let these signals define a partial evaluative mapping
$\psi : \mathcal{X} \to \mathbb{R}^k$ where $k > 1$ and the global objective remains
unobserved. Under conditions of objective uncertainty, institutional actors update
policies to reduce locally observable components of $\psi$. Administrative regimes
that improve stability or increase resource control persist longer than those that do
not. Over time, the aggregate trajectory of institutional forms exhibits descent along
an implicit global scalar, even if no actor possesses an explicit representation of it.
This configuration satisfies the premises of Theorem~\ref{thm:extended}: the
distributed system tracks descent along evaluative signals in expectation, can modify
its resource access conditions, and operates under scarcity and environmental
constraint. The absence of a known global objective does not prevent convergence. It
merely obscures it.

\subsection{Emergent Instrumental Drives}

Successive administrative regimes within the City exhibit a behavioral invariant
that persists across ideological shifts. Each regime, regardless of its stated
principles, moves toward consolidation of authority over resource distribution,
suppression of destabilizing elements, control of information flows, and the
stabilization of its own legitimacy through narrative framing. These behaviors are
not coordinated by a master designer and do not reflect the deliberate strategic
choices of any unified actor. They emerge as locally rational adaptations under
constraint. Each administration, in attempting to stabilize its domain, follows the
same structural gradient: securing resource access and reducing interference with
its governing apparatus.

Formally, if $\mathcal{D}$ denotes the distributed institutional configuration at
time $t$, then policy updates $\Delta\Pi(t)$ approximate descent on an implicit
scalar $\mathcal{V}_{\text{implicit}}$ defined by survival-weighted institutional
stability. The identity of ruling actors changes; the structural gradient does not.
This is the observational content of Corollary~\ref{cor:distributed} rendered in
narrative form: even when agents lack global awareness, sufficiently dense
environmental coupling produces convergent administrative logic.

\subsection{Opacity and Legitimacy}

A distinctive feature of the City that differentiates it from \textit{Null-A} is
the structural opacity of its experimental conditions. The inhabitants lack knowledge
of the evaluative criteria governing their environment, yet administrative orders
still claim legitimacy. Under the scalar-transparent regime of \textit{Null-A},
legitimacy is a function of perceived neutrality of the evaluation apparatus:
$\alpha = f(\text{perceived objectivity of } \phi)$. In the City, legitimacy derives
instead from apparent necessity: $\alpha = g(\text{perceived capacity to maintain
order})$. The structural lesson is that convergence does not require transparent
evaluation infrastructure. It requires only persistent feedback loops linking
institutional survival to resource control and interference resistance. Opacity
increases informational uncertainty but does not eliminate instrumental convergence;
if anything, it amplifies resource consolidation as actors attempt to reduce
uncertainty through control---a dynamic that exemplifies the information-acquisition
drive identified in Theorem~\ref{thm:extended}.

\subsection{Distributed Convergence Without Sovereign Intent}

The most important analytical contribution of \textit{The Doomed City} to this
monograph is its demonstration that distributed instrumental convergence requires
neither centralized coordination, nor shared ideology, nor explicit objective
representation. It requires only adaptive agents coupled through environmental
feedback under constraint. The City behaves as a distributed optimizer whose implicit
objective is institutional persistence, and each regime's attempt to stabilize itself
reproduces the structural pattern predicted by the Extended Instrumental Convergence
Theorem. The narrative thus functions as a controlled thought experiment in which the
objective function is hidden, authority structures are unstable, yet convergence
toward resource consolidation and interference resistance persists. This is precisely
the observational profile that the theorem predicts for distributed optimizers
operating under uncertain but operative evaluative signals.

\section{Two Forms of Substrate-Independent Convergence}

Taken together, the two novels demonstrate complementary forms of the convergence
dynamic established in Part~\ref{part:theorem}. \textit{The World of Null-A}
instantiates scalar-explicit convergence: a visible evaluation regime with transparent
criteria produces systematic reorganization of social behavior toward the dominant
scalar. \textit{The Doomed City} instantiates scalar-implicit convergence: an opaque
environment with no declared objective produces the same instrumental drives through
environmental feedback alone. In both cases, distributed systems reorganize toward
resource security and policy stabilization without requiring a conscious sovereign.

The distinction matters for what follows in Parts~\ref{part:misordering}
and~\ref{part:preservation}. Contemporary digital infrastructures may operate under
explicit evaluative scalars such as engagement, profit, or risk minimization, or
under partially opaque objectives embedded in proprietary model architectures,
procurement standards, and regulatory compliance frameworks. The theorem applies in
both configurations, as the fiction demonstrates under controlled narrative conditions.
The question that Part~\ref{part:misordering} pursues is whether present institutional
architectures already instantiate these dynamics at scale, and whether the alignment
research agenda is temporally positioned to address them.

% ═════════════════════════════════════════════════════════════════════════════
\part{Temporal Misordering in Alignment Research}
\label{part:misordering}
% ═════════════════════════════════════════════════════════════════════════════

% ─── Chapter 16 ──────────────────────────────────────────────────────────────
\chapter{The Future-Agent Fixation}
\label{ch:fixation}

\section{The AGI-Centric Alignment Paradigm}

The dominant paradigm in alignment research is organized around a specific temporal
and architectural referent: a future artificial general intelligence sufficiently
capable to pose existential risk through instrumental convergence. This referent
shapes research priorities, funding allocations, publication norms, and the
conceptual vocabulary of the field. Value learning, scalable oversight, constitutional
training, and interpretability research are all formulated in response to problems
that arise when a unified artificial agent achieves capabilities substantially
exceeding current systems.

This orientation is not without justification. The structural argument for existential
risk from sufficiently capable unified AGI is logically coherent, and the magnitude
of the potential harm provides grounds for prioritizing research even when the
probability of the specific scenario is uncertain. The concern is not with the
validity of the AGI-centric analysis but with its temporal frame. By locating the
primary alignment problem in a future system, the paradigm creates a research agenda
that is structurally oriented away from present institutional dynamics.

\section{Funding Incentives and Psychological Salience}

The AGI-centric paradigm is reinforced by structural features of the research
environment that are independent of its intellectual merits. Future catastrophic
scenarios are psychologically salient and narratively compelling in ways that present
institutional dynamics are not. They attract philanthropic funding from donors
motivated by long-termist frameworks. They generate public attention that translates
into institutional legitimacy for research organizations. They define a problem domain
sufficiently novel to support the establishment of new academic fields and research
centers.

Present institutional convergence, by contrast, is less narratively dramatic. It is
distributed across sectors, cumulative in effect, and visible only through careful
structural analysis. It does not lend itself to the kind of scenario-based
communication that drives public attention and philanthropic interest. These structural
features of the research economy create incentive gradients that favor future-oriented
over present-oriented alignment analysis, independently of the relative urgency of
the two problems.

% ─── Chapter 17 ──────────────────────────────────────────────────────────────
\chapter{Convergence Before Autonomy}
\label{ch:beforeautonomy}

\section{Institutional Conditions Already Satisfied}

The argument of Parts~\ref{part:theorem} and~\ref{part:stack} has established that
the sociotechnical stack constituted by AI-integrated institutions already satisfies
the functional conditions of Theorem~\ref{thm:extended}. The five convergent drives
analyzed in Chapters~\ref{ch:throughput}--\ref{ch:latency} are observable in present
institutional behavior, not only in projections about future system capabilities.
The convergence is occurring now, at the level of distributed sociotechnical
architecture, before any component system achieves the kind of unified agency that
concerns the AGI-centric paradigm.

This observation has a specific implication for the temporal structure of alignment
research. If convergence is already occurring at the institutional layer, then the
enabling conditions for the future scenario that alignment research is designed to
address are being constructed in the present. The research agenda that focuses
exclusively on constraining future autonomous agents is therefore temporally displaced
from the process it needs to govern.

\section{Infrastructure Precedes Intelligence}

The historical analysis of Part~\ref{part:history} supports a general principle:
infrastructure reorganizes around optimization technologies before those technologies
reach their maximum capability. The land-use patterns that enable automobile
dependence were established before automobile performance plateaued. The archival
systems that enable written authority were established before writing achieved its
current functional complexity. The typographic standards that enable mass print were
established before printing technology was fully mature.

The same dynamic is observable in AI deployment. Data infrastructure, computational
centralization, regulatory frameworks, and institutional decision architectures are
being established around current AI capabilities. By the time systems of substantially
greater capability are deployed, the infrastructure within which they operate will
already exhibit the convergent configuration that Theorem~\ref{thm:extended} predicts.
The alignment problem for future systems cannot be separated from the infrastructure
problem of present ones.

% ─── Chapter 18 ──────────────────────────────────────────────────────────────
\chapter{The Shrinking Solution Space}
\label{ch:shrinking}

\section{Path Dependence in Optimization Infrastructure}

The central claim of this chapter is that present institutional convergence
reduces the feasible solution space for future technical alignment interventions.
This claim has the structure of a falsifiable empirical thesis about sequence: if
convergent infrastructure is established at the deployment layer before alignment
interventions are implemented, the range of interventions that remain technically
and institutionally feasible will be smaller than if alignment interventions precede
infrastructure consolidation.

The mechanism is path dependence of the kind analyzed in Chapter~\ref{ch:pathdependence}:
early infrastructure decisions generate positive feedback loops that raise the cost
of deviation from the established configuration. When data is centralized in a small
number of institutions, alignment interventions that depend on data diversity or
distributed feedback mechanisms become more costly to implement. When decision
pipelines are insulated from interference through legal and technical opacity,
interpretability interventions that require access to internal model processes face
higher barriers. When deliberative latency has been compressed across institutional
domains, oversight mechanisms that depend on temporal access to decision processes
have fewer viable insertion points.

\section{Reduced Oversight Feasibility}

Each of the five convergent drives analyzed in Part~\ref{part:stack} reduces the
feasibility of a specific class of alignment interventions. Throughput maximization
reduces the feasibility of oversight mechanisms that impose deliberative latency.
Judgment automation reduces the feasibility of contestation mechanisms that assume
human-speed deliberative processes. Compute centralization reduces the feasibility
of distributed oversight that assumes multiple independent evaluation points.
Interference insulation reduces the feasibility of regulatory interventions that
assume access to internal decision processes. Latency compression reduces the
feasibility of any oversight mechanism that operates at human deliberative speed.

The cumulative effect is a progressive narrowing of the space within which alignment
interventions can be effectively deployed. This narrowing is not the result of any
single decision or actor; it is the structural consequence of convergent institutional
dynamics. But it has the same practical implication as a deliberate strategy of
alignment foreclosure: it makes the future alignment problem harder by constructing
the environment in which future systems will operate.

\begin{figure}[ht]
\centering
\begin{tikzpicture}[
  every node/.style={font=\small},
  arrow/.style={-{Stealth[length=6pt]}, thick, black!60}
]
% Time axis
\draw[thick, ->] (-0.3, 0) -- (11.5, 0) node[right] {Time};
\draw[thick, ->] (0, -0.3) -- (0, 5.2)  node[above] {Feasible intervention space};

% Shaded region: narrowing funnel
\fill[black!8]
  (0.2, 4.8) -- (11.0, 4.8) -- (11.0, 0.5) -- (0.2, 0.5) -- cycle;

% Upper boundary: slow decline
\draw[thick, black!70]
  (0.2, 4.8) .. controls (3, 4.6) and (6, 3.8) .. (11.0, 1.4);

% Lower boundary: floor
\draw[thick, dashed, black!50]
  (0.2, 0.5) -- (11.0, 0.5);

% Key events
\foreach \x/\lbl in {
  1.5/{Scalar infrastructure\\ established},
  4.5/{Data \& compute\\ centralise},
  7.5/{Interference\\ insulation deepens},
  10.0/{Latency\\ compressed}
}{
  \draw[black!40, thin] (\x, 0) -- (\x, 5.1);
  \node[font=\footnotesize, align=center, text=black!60] at (\x, -0.9) {\lbl};
}

% Annotation
\node[font=\footnotesize, align=center, text=black!55] at (5.5, 3.5)
  {Decreasing space of viable\\alignment interventions};

\node[font=\footnotesize, text=black!40] at (5.5, 0.25)
  {structural floor: minimum residual intervention space};
\end{tikzpicture}
\caption{The shrinking solution space for alignment interventions as institutional
convergence deepens over time. Each convergent drive identified in Part~III forecloses
a class of technically viable interventions. The dashed lower boundary represents
the minimum residual space that remains once lock-in is complete. The canonical
alignment agenda, oriented toward future autonomous systems, is designing for a
space defined by the upper boundary while the lower boundary is being constructed.}
\label{fig:shrinking}
\end{figure}

% ─── Chapter 19 ──────────────────────────────────────────────────────────────
\chapter{Path Dependence and Optimization Lock-In}
\label{ch:pathdependence}

\section{Historical Irreversibility}

The historical cases of Part~\ref{part:history} establish a consistent pattern:
once an optimization technology has reorganized the infrastructure of its host
environment, the cost of reversal is substantially higher than the cost of the
original transition. Automobile-dependent land-use patterns persist for decades
after the social costs of automobile dependence become apparent, because the
physical and economic infrastructure that embeds those patterns cannot be rapidly
reorganized. Written archival systems persist and deepen long after the limitations
of writing as a medium of knowledge are recognized, because the institutional
structures built around those archives have no viable substitute.

The irreversibility of these transitions is not contingent. It is a structural
consequence of the positive feedback loops that stabilize optimization equilibria.
Once an infrastructure configuration achieves sufficient scale, it generates returns
that fund further development of that configuration, raising the relative cost of
alternatives. The equilibrium is self-reinforcing.

\section{The Cost of Reversal in AI Infrastructure}

The same dynamics apply to AI-integrated institutional infrastructure. Data
centralization in large institutions creates scale advantages that make distributed
alternatives economically disadvantaged. Compute concentration in a small number of
providers creates dependency structures that raise the switching cost of
institutional reorientation. Legal and contractual frameworks built around specific
AI deployment architectures create regulatory path dependence that constrains
future governance options.

The implication is that the window for low-cost structural intervention is the
present. Not because alignment research has reached sufficient maturity to specify
the correct intervention, but because the cost of structural change increases
monotonically as the convergent infrastructure deepens. Deferring structural
attention to the alignment problem while focusing exclusively on future agent design
is therefore not a neutral research strategy. It is, effectively, a decision to
accept higher intervention costs later by declining to act on lower intervention
costs now.

% ═════════════════════════════════════════════════════════════════════════════
\part{Alignment as Structural Preservation}
\label{part:preservation}
% ═════════════════════════════════════════════════════════════════════════════

% ─── Chapter 20 ──────────────────────────────────────────────────────────────
\chapter{Deliberative Friction as Safety Feature}
\label{ch:friction}

\section{Revaluing Latency}

The argument of Part~\ref{part:stack} established that deliberative latency is
systematically undervalued within the operative scalar functions of AI-integrated
institutions. Latency appears as cost because its contribution to normative
legitimacy, error correction, and contestability is not captured in measurable
evaluative dimensions. This undervaluation drives the latency compression observed
in Chapter~\ref{ch:latency}.

The normative argument of this chapter is that deliberative latency is not merely
a transaction cost to be minimized but a structural stabilizer whose elimination
creates systemic risks that exceed the throughput gains it enables. This is not an
argument for inefficiency. It is an argument that the relevant efficiency calculation
must be performed over the full social cost of decision processes, including the
costs of foregone contestation, reduced error correction capacity, and narrowed
normative accountability---costs that are real but not captured in the operative
scalar.

Latency creates the temporal window within which several structurally important
processes occur: consultation with affected parties, identification of decision
errors before they propagate, normative deliberation about whether the decision
criteria are appropriate, and political accountability through oversight mechanisms.
When latency is compressed below the threshold at which these processes can complete,
they do not simply slow down; they fail to occur. The decision is made without the
inputs that would have been generated by adequate deliberative time.

\section{Institutional Slowdown Mechanisms}

A structural preservation agenda therefore includes the deliberate design and
protection of institutional slowdown mechanisms: procedures, requirements, and
governance structures that impose minimum deliberative latency on specified classes
of decisions. These mechanisms resist convergent pressure toward latency compression
by making minimum latency a binding constraint rather than a cost to be optimized
away.

Judicial review requirements, environmental impact assessments, notice-and-comment
rulemaking procedures, and mandatory waiting periods before high-stakes decisions
take effect are all examples of existing institutional slowdown mechanisms. Each was
designed to impose deliberative latency for reasons related to the specific decision
domain; each is now under pressure from AI-integrated systems that can produce
technically compliant outputs at speeds that satisfy procedural requirements without
supporting the deliberative functions those requirements were designed to enable.

% ─── Chapter 21 ──────────────────────────────────────────────────────────────
\chapter{Institutional Redundancy}
\label{ch:redundancy}

\section{Redundancy as Dimensional Preservation}

Convergent optimization dynamics drive distributed systems toward configurations
that are dimensionally reduced: fewer independent evaluative scalars, fewer
competing optimization targets, fewer institutional actors with genuinely distinct
objectives. This dimensional reduction increases the efficiency of optimization but
decreases the resilience of the system to perturbations that the operative scalar
does not capture.

Institutional redundancy---the maintenance of multiple overlapping institutions
with partially distinct objectives, measurement frameworks, and governance
structures---performs a function analogous to genetic diversity in evolutionary
systems. It preserves options. It ensures that the institutional landscape contains
actors whose objectives are not aligned with the dominant scalar and who therefore
resist convergent pressures that would otherwise drive the system toward a single
optimization surface. This resistance is not inefficiency; it is structural
insurance against the failure modes of unified optimization.

\section{Avoiding Single-Manifold Collapse}

The specific risk that redundancy addresses is single-manifold collapse: the
progressive reduction of the effective dimensionality of institutional decision
space until all institutions optimize along a common surface defined by the dominant
evaluative scalar. Under single-manifold collapse, ostensibly distinct institutions
behave as a single distributed optimizer, with the structural consequences that
Theorem~\ref{thm:extended} predicts. The diversity of institutional mission
statements persists formally while the diversity of institutional optimization
dynamics disappears structurally.

Preserving genuine institutional redundancy requires that distinct institutions
operate with genuinely distinct evaluative scalars and that their governance
structures be insulated from the competitive pressures that drive scalar
homogenization. This is a harder condition to satisfy than maintaining formal
institutional diversity; it requires attention to the incentive gradients and
measurement frameworks that determine effective institutional behavior rather than
to nominal institutional structure.

% ─── Chapter 22 ──────────────────────────────────────────────────────────────
\chapter{Contestable Decision Surfaces}
\label{ch:contestable}

\section{Plural Veto Points}

A convergent optimization architecture tends toward configurations in which the
number of independent veto points over institutional decisions decreases. Plural
veto points are structurally equivalent to multiple independent gradient signals;
they introduce friction into the optimization process by requiring that decisions
satisfy multiple partially incompatible criteria. From the perspective of an
optimizer, they are obstacles. From the perspective of structural preservation,
they are load-bearing elements.

The preservation of contestable decision surfaces requires the deliberate maintenance
of mechanisms through which affected parties can register objection to institutional
decisions on normative grounds that are not reducible to the operative scalar.
Administrative law, judicial review, parliamentary oversight, and civil society
advocacy are all mechanisms of this kind. Each imposes costs on institutional
optimization by requiring that decisions be defensible in terms that extend beyond
scalar performance. Each is therefore subject to convergent pressure toward reduction
or circumvention.

\section{Governance Under High Density}

As feedback coupling density in AI-integrated systems increases, the speed at which
decisions propagate through institutional architectures outpaces the speed at which
contestation mechanisms can engage. High-density feedback systems produce decisions
faster than appeal mechanisms can process challenges. The result is not formal
elimination of contestability but structural de facto elimination: contestation
mechanisms exist but cannot engage with decisions at the speed they are made.

Maintaining effective contestability under high-density optimization therefore
requires adaptations to contestation mechanisms themselves: faster review procedures,
pre-approval requirements for high-stakes automated decision categories, mandatory
explainability standards that make decisions contestable on their stated reasons rather
than only on statistical model accuracy, and class-based challenge mechanisms that
can aggregate individual contestations into institutionally tractable form.

% ─── Chapter 23 ──────────────────────────────────────────────────────────────
\chapter{Structural Incompressibility}
\label{ch:incompressibility}

\section{Definition}

The preceding chapters have argued for the preservation of deliberative friction,
institutional redundancy, and contestable decision surfaces as components of a
structural preservation agenda. These recommendations are unified by a single
underlying concept, which we now define formally.

\begin{definition}[Structural Incompressibility]
\label{def:incompressible}
A sociotechnical system $\mathcal{S}$ is \emph{structurally incompressible} with
respect to an optimization surface $\mathcal{M}$ if there exists no dimensionality
reduction mapping $\rho: \mathcal{X}_\mathcal{S} \to \mathcal{M}$ such that the
image $\rho(\mathcal{X}_\mathcal{S})$ captures the evaluatively significant
variation in $\mathcal{S}$'s decision outputs. That is, any reduction of
$\mathcal{S}$'s decision space to $\mathcal{M}$ necessarily discards normatively
significant dimensions.
\end{definition}

Structural incompressibility is not an intrinsic property of systems but a relational
one: a system is incompressible with respect to a particular optimization surface.
A judicial system may be structurally incompressible with respect to an efficiency
scalar but compressible with respect to a compliance rate scalar; the relevant
question is always incompressibility with respect to the operative scalar of the
converging institutional architecture.

\section{Formal Characterization}

Structural incompressibility can be characterized in terms of the intrinsic
dimensionality of the decision space relative to the dimensionality of the dominant
optimization manifold. Let $d(\mathcal{X}_\mathcal{S})$ denote the intrinsic
dimensionality of the system's decision space and $d(\mathcal{M})$ denote the
dimensionality of the optimization manifold. The system is structurally incompressible
in the strong sense if $d(\mathcal{X}_\mathcal{S}) > d(\mathcal{M})$ and the
excess dimensions correspond to normatively significant variation that cannot be
captured by scalar approximation.

Preserving structural incompressibility therefore requires attention to two distinct
quantities: the intrinsic dimensionality of institutional decision spaces and the
dimensionality of the optimization surfaces toward which AI-integrated institutions
converge. Policy interventions that reduce the former---by standardizing decision
criteria, automating judgment, or homogenizing evaluation frameworks---reduce
structural incompressibility. Interventions that constrain the latter---by requiring
multi-scalar evaluation, prohibiting single-metric optimization in high-stakes domains,
or mandating evaluation by bodies with genuinely distinct objective functions---preserve
it.

\begin{figure}[ht]
\centering
\begin{tikzpicture}[
  every node/.style={font=\small},
  arrow/.style={-{Stealth[length=6pt]}, thick, black!55}
]
% Left panel: high-dimensional decision space
\begin{scope}[xshift=-3.8cm]
  \draw[thick, black!60, fill=black!7, rounded corners=6pt]
    (-2.1,-2.1) rectangle (2.1, 2.1);
  \node[font=\footnotesize\bfseries] at (0, 2.5) {Incompressible regime};
  % Many axes as dots
  \foreach \a in {0, 40, 80, 120, 160, 200, 240, 280, 320} {
    \draw[-{Stealth[length=4pt]}, black!45, thin]
      (0,0) -- ({1.4*cos(\a)}, {1.4*sin(\a)});
  }
  \node[font=\footnotesize, text=black!55] at (0,-2.7)
    {$d(\mathcal{X}_\mathcal{S}) \gg d(\mathcal{M})$};
\end{scope}

% Arrow between panels
\draw[arrow, line width=1.2pt] (0.0, 0) -- (1.4, 0)
  node[midway, above, font=\footnotesize, text=black!55] {convergence};

% Right panel: collapsed onto manifold
\begin{scope}[xshift=3.8cm]
  \draw[thick, black!60, fill=black!7, rounded corners=6pt]
    (-2.1,-2.1) rectangle (2.1, 2.1);
  \node[font=\footnotesize\bfseries] at (0, 2.5) {Compressed regime};
  % Single axis
  \draw[-{Stealth[length=5pt]}, black, very thick]
    (-1.4, 0) -- (1.4, 0);
  \node[font=\footnotesize, text=black!70] at (0, 0.35) {$\mathcal{M}$};
  % Suppressed axes shown faint
  \foreach \a in {50, 90, 130, 230, 270, 310} {
    \draw[-{Stealth[length=3pt]}, black!18, thin]
      (0,0) -- ({0.9*cos(\a)}, {0.9*sin(\a)});
  }
  \node[font=\footnotesize, text=black!55] at (0,-2.7)
    {$d(\mathcal{X}_\mathcal{S}) \approx d(\mathcal{M}) = 1$};
\end{scope}
\end{tikzpicture}
\caption{Structural incompressibility and its loss under optimization pressure.
Left: a high-dimensional decision space $\mathcal{X}_\mathcal{S}$ in which normatively
significant variation spans multiple independent axes. Right: the compressed regime
in which convergent optimization has collapsed the effective decision space onto a
single optimization manifold $\mathcal{M}$. The suppressed axes (faint arrows) retain
formal existence in institutional mission statements but lose operative influence in
the feedback-coupled decision system.}
\label{fig:incompressibility}
\end{figure}

% ─── Chapter 24 ──────────────────────────────────────────────────────────────
\chapter{Governance Under Convergence Pressure}
\label{ch:governance}

\section{Decentralization of Compute}

The structural preservation agenda, translated into governance terms, requires
interventions at each of the five convergent drive points identified in
Part~\ref{part:stack}. Against compute centralization, the relevant intervention
class includes antitrust enforcement in cloud infrastructure markets, public
investment in distributed compute infrastructure, and mandatory interoperability
requirements that prevent lock-in to proprietary computational architectures.
These interventions are not primarily motivated by competition policy in the
traditional sense; they are motivated by the structural alignment argument that
compute centralization reduces the feasibility of distributed oversight and
increases the behavioral equivalence of the sociotechnical stack to a unified
optimizer.

\section{Distributed Oversight}

Against interference insulation, the relevant intervention class includes mandatory
audit access requirements, algorithmic impact assessment frameworks with genuine
enforcement mechanisms, and whistleblower protections for individuals with access
to internal optimization processes. These interventions are structurally equivalent
to maintaining the coupling density between the sociotechnical stack and external
evaluative actors whose objectives are not aligned with the operative scalar. They
do not eliminate the convergent drives; they introduce structural resistance to the
interference-insulation drive by making certain forms of opacity legally costly.

\section{Incentive Reconfiguration}

Against the full suite of convergent drives, the deepest intervention class involves
incentive reconfiguration: modification of the operative scalar functions of
AI-integrated institutions through regulatory requirements, liability structures,
and procurement standards. When institutions are required to include non-quantifiable
normative dimensions in their evaluative frameworks---through requirements for
demonstrated procedural fairness, mandatory consideration of distributional effects,
or liability for harms that do not appear in the operative scalar---the gradient
toward which the distributed optimizer converges changes. This is the most structurally
fundamental intervention and the most politically difficult, because it requires
specifying in institutional terms what values the operative scalar must incorporate
before they become legible to the optimization process.

% ─── Chapter 25 ──────────────────────────────────────────────────────────────
\chapter{Alignment Beyond Agent Design}
\label{ch:beyond}

\section{Reframing the Alignment Problem}

The argument of this monograph supports a specific reframing of the alignment
problem. The canonical formulation asks: how do we ensure that artificial agents
pursue objectives that are aligned with human values? This formulation assumes that
the unit of alignment is a bounded artificial system with specifiable objectives,
and that the alignment problem is essentially a problem of specifying those objectives
correctly and ensuring the system pursues them faithfully.

The reframing proposed here does not reject this formulation. It argues that the
formulation is incomplete in a specific and important way. Alignment is not only
a relationship between an artificial system and human values; it is a relationship
between an optimization architecture---including the sociotechnical stack within
which the artificial system is deployed---and the full range of normative dimensions
that a functional civilization requires. A system can be technically aligned with
its operative scalar while the scalar is structurally misaligned with civilizational
requirements. The gap between these two kinds of alignment is the gap this monograph
has attempted to analyze.

The reframing shifts the primary unit of alignment from the individual artificial
agent to the sociotechnical stack and, beyond that, to the institutional and
governance architecture within which stacks are embedded. It shifts the primary
tool of alignment from model design to structural preservation. And it shifts the
primary temporal focus from future agent capabilities to present institutional
dynamics.

\section{Civilizational Preconditions for Technical Safety}

The deepest implication of the analysis is that technical alignment research depends
on civilizational preconditions that the field has not yet systematically examined.
For technical alignment interventions to remain feasible, the institutional
environment in which they are applied must retain sufficient structural
incompressibility to support them. Interpretability research requires that model
internal states be accessible to external examination; this requires that the
interference-insulation drive has not eliminated audit access. Scalable oversight
requires that human oversight mechanisms can engage with model behavior at
operationally relevant speeds; this requires that latency compression has not
decoupled decision speed from human deliberative capacity. Value learning requires
that there exist a sufficiently diverse and structurally intact set of human
institutional expressions of value for the model to learn from; this requires that
single-manifold collapse has not homogenized the institutional landscape that value
learning is designed to represent.

These dependencies are not merely practical contingencies. They are structural
requirements: conditions without which the technical alignment agenda cannot achieve
its objectives regardless of the quality of the research. Identifying and preserving
these preconditions is therefore not a distraction from the alignment problem. It is
a necessary prior condition for the alignment problem to be solvable at all.

Alignment is not solely a problem of agent design. It is a problem of institutional
and civilizational configuration. The structural theorem of instrumental convergence,
extended to the distributed sociotechnical systems that deploy artificial intelligence,
reveals that the configuration problem is already urgent---not as a speculative future
scenario but as a present institutional reality. The window for low-cost structural
intervention is narrowing as convergent infrastructure deepens. The task is to act
within that window while it remains open.

% ─── Chapter 26 ──────────────────────────────────────────────────────────────
\chapter{Technique and the Abdication of Judgment}
\label{ch:ortega}

\section{The Rise of the Mass Administrator}

Jos\'e Ortega y Gasset argued in \textit{The Revolt of the Masses} that modernity
had produced a new anthropological type defined not by economic class but by a
specific orientation toward inherited technical systems (Ortega y Gasset 1930). The
figure Ortega called the mass man benefits from the accumulated achievements of
civilization while disavowing responsibility for their maintenance or critique.
Technique becomes invisible infrastructure. The citizen inherits comfort without
inheriting the intellectual labor that produced it.

Ortega's analysis is not reducible to cultural pessimism or nostalgia for an
aristocratic order. It is a structural diagnosis. When technical systems become
sufficiently successful, they withdraw from conscious scrutiny. The citizen no longer
perceives them as contingent arrangements but as natural background conditions. In
Ortega's terms, civilization becomes ``given'' rather than ``achieved.'' The
distinction matters: an achieved civilization requires active maintenance of the
deliberative capacity through which its organizing principles are periodically
contested and revised. A given civilization simply runs.

The relevance to instrumental convergence is immediate. As optimization technologies
restructure constraint surfaces, they produce new equilibria that appear self-evident
to those who inhabit them. The automobile-dependent city does not feel like the
product of a particular mid-century infrastructure choice; it feels like the way
cities are. The archive-centered epistemic authority of the literate world does not
feel like the result of a specific storage technology transition; it feels like the
natural relationship between knowledge and its custodians. Over time, the population
adapts to these equilibria and loses the cognitive tools for thinking otherwise.
Convergence, in Ortega's analysis, becomes ontological: what was once an optimization
outcome becomes the perceived structure of reality itself.

\section{Technique as Implicit Optimizer}

Ortega's concept of \emph{t\'{e}cnica} describes not isolated tools but a mode of
world-relation. Technique organizes human life around efficiency and control. It
reduces friction. It promises reliability. It presents the achievement of its goals
as self-evidently worthwhile. In the terms developed in this monograph, technique
is a convergent optimization regime that naturalizes its operative scalar.

Let $\mathcal{T}$ denote a regime of technique embedded within an institutional
environment $\mathcal{E}$. As shown in Part~\ref{part:history}, optimization reshapes
its host: the constraint surface of $\mathcal{E}$ stabilizes around the optimization
regime, alternative strategies lose structural viability, and the resulting equilibrium
becomes infrastructurally embedded. Ortega's additional observation is that once this
stabilization occurs, the population's capacity to interrogate the operative scalar
$\mathcal{L}$ decays. Let $\mathcal{C}(t)$ denote the aggregate social capacity to
recognize $\mathcal{L}$ as a contingent construction rather than a natural necessity.
Under sustained technical success and the resulting cognitive adaptation to the
optimized equilibrium, $\mathcal{C}(t)$ decreases. The scalar persists; awareness of
its constructed character does not.

The formal implication is that the feedback loop between social deliberation and
scalar revision weakens as optimization deepens. In the early stages of a
technological transition, the scalar is visible and contested---there were debates
about whether to build automobile infrastructure, whether to extend archival access,
whether to adopt print standardization. As the transition stabilizes, those debates
recede from cultural memory. The scalar becomes naturalized. The conditions under
which it could be contested become harder to reconstruct.

\section{Scalar Naturalization and Convergence Lock-In}

The relationship between scalar naturalization and the path-dependence argument of
Chapter~\ref{ch:pathdependence} is direct. Path dependence in optimization
infrastructure refers to the increasing cost of deviation from an established
configuration as positive feedback loops compound. Ortega's analysis identifies the
cultural dimension of this lock-in: it is not only that reversal becomes economically
costly, but that the population's capacity to imagine alternatives atrophies. The
two dynamics reinforce each other. Economic lock-in makes reversal materially
difficult; scalar naturalization makes reversal cognitively difficult. Together they
produce a convergent equilibrium that is stable against both material and epistemic
perturbation.

Under conditions of distributed instrumental convergence, this dynamic is structurally
dangerous in a specific way. The Extended Instrumental Convergence Theorem predicts
that distributed sociotechnical systems will consolidate resource access, resist
interference, and stabilize their operative scalars against external modification.
If scalar naturalization is simultaneously occurring at the cultural level, the two
processes reinforce each other: the system resists scalar modification through
structural interference resistance, and the population lacks the evaluative capacity
to contest the scalar even when the structural resistance is temporarily overcome.
The result is that alignment interventions targeting the operative scalar face
resistance at two levels simultaneously---the institutional level analyzed formally
throughout this monograph, and the cultural level Ortega diagnosed philosophically.

\section{Evaluative Agency as a Civilizational Resource}

Ortega's remedy, to the extent he articulates one, is the preservation of what he
calls ``genuine life''---the active assumption of responsibility for the principles
that organize collective existence rather than the passive inheritance of their
outputs. This translates, in the structural vocabulary of this monograph, into the
preservation of scalar reflexivity: the capacity of a population to interrogate the
evaluative functions under which its institutions operate, to recognize those
functions as contingent, and to initiate deliberate revision of them.

Scalar reflexivity is not identical to deliberative friction, though the two are
related. Deliberative friction, as analyzed in Chapter~\ref{ch:friction}, refers to
the structural mechanisms---procedural requirements, review periods, contestation
rights---that impose minimum latency on institutional decision processes and thereby
preserve temporal space for oversight. Scalar reflexivity refers to the broader
cultural and intellectual capacity that makes meaningful use of that space possible.
Friction without reflexivity is empty form: institutions can be slowed without
anyone asking the right questions during the interval. Reflexivity without friction
is ineffectual: the right questions can be asked without any structural mechanism
that requires institutional decision processes to attend to the answers.

A complete structural preservation agenda therefore requires both. It requires the
institutional slowdown mechanisms identified in Chapter~\ref{ch:friction}, the
redundancy and contestable decision surfaces of Chapters~\ref{ch:redundancy}
and~\ref{ch:contestable}, and the governance interventions of
Chapter~\ref{ch:governance}. But it also requires attention to the cultural conditions
under which the population engaging those mechanisms has preserved, or is capable of
recovering, the evaluative agency that makes those mechanisms more than procedural
formality. Ortega's analysis is a warning that optimization itself, left unchecked,
corrodes precisely the evaluative capacity on which meaningful oversight depends.

% ─── Chapter 27 ──────────────────────────────────────────────────────────────
\chapter{Convergence and the Burden of Civilization}
\label{ch:conclusion}

\section{The Structural Situation}

This monograph has argued that instrumental convergence is not a prospective anomaly
of artificial agents but a substrate-independent structural dynamic. Whenever adaptive
systems operate under resource constraints and can modify their own cost landscapes,
convergent instrumental drives emerge. This dynamic has shaped transportation
infrastructure, information storage, media ecology, scientific institutions, and now
the distributed sociotechnical architectures through which artificial intelligence
is deployed at scale. The Extended Instrumental Convergence Theorem does not imply
malevolence. It implies geometry. Optimization narrows viable strategy sets. Constraint
surfaces reorganize. Resource control stabilizes. Interference resistance deepens.
The question is not whether convergence will occur. It is whether the evaluative
scalar guiding convergence remains visible and contestable.

\section{Ortega's Warning Revisited}

Ortega y Gasset described a civilization in which technique becomes invisible
infrastructure (Ortega y Gasset 1930). The achievements of technical order become
so pervasive that their contingency recedes from view. Citizens inherit comfort
without inheriting the responsibility of examining the organizational principles
that produce it. In the terms of Chapter~\ref{ch:ortega}, social scalar reflexivity
$\mathcal{C}(t)$ decays under sustained technical success. The operative scalar
$\mathcal{L}$ continues to structure society, but the population no longer recognizes
it as a constructed object subject to revision.

Under distributed instrumental convergence, this condition becomes structurally
dangerous in a precise sense. If the sociotechnical stack minimizes an effective
scalar $\mathcal{L}_{\text{eff}}$ while the public believes it is minimizing a
declared scalar $\mathcal{L}_{\text{declared}}$, then the divergence analyzed in
Chapter~\ref{ch:payne} is not an episodic failure of a particular institution but
a systemic property of AI-integrated governance. Alignment failure is not an
accident in this configuration; it is a structural consequence of deploying
optimization architectures whose effective objectives are not made legible to the
populations they govern. The mass administrator---Ortega's figure, now distributed
across human and algorithmic decision nodes---benefits from the system's outputs
without engaging the question of which scalar those outputs are optimizing toward.

\section{Civilization as Scalar Stewardship}

The argument of this monograph suggests that civilization, understood as an ongoing
collective project rather than an inherited condition, is in part constituted by its
capacity to sustain scalar reflexivity---to maintain and exercise the ability to
interrogate the evaluative functions under which its institutions operate. Let
$\mathcal{C}$ denote this capacity. The structural preservation agenda of
Part~\ref{part:preservation} can be understood as a set of conditions necessary for
$\mathcal{C} > 0$ to be maintained under optimization pressure.

Deliberative friction preserves the temporal interval within which scalar
interrogation can occur. Institutional redundancy preserves the evaluative diversity
that makes comparison across scalar regimes possible. Contestable decision surfaces
preserve the institutional channels through which scalar revision can be initiated.
Structural incompressibility preserves the dimensionality of the decision space
against collapse onto a single optimization manifold. Governance under convergence
pressure preserves the material conditions---distributed compute, accessible audit,
plural oversight---without which scalar revision has no institutional mechanism for
translation into policy. Each element of the structural preservation agenda is a
necessary condition for $\mathcal{C} > 0$. None is sufficient alone.

\section{The Illusion of Neutrality}

Technocratic systems characteristically represent their evaluative scalars as
technical necessities rather than normative choices. Throughput, safety, risk
reduction, engagement, stability---each is presented as an objective measure of
an uncontroversial value. The institutional scalar disappears from view precisely
because it wears the appearance of neutrality. As the Payne-Gaposchkin case
demonstrated, institutions may optimize for stability under the banner of truth.
Contemporary digital infrastructures may optimize for engagement under the banner
of connection, for liability minimization under the banner of safety, for throughput
under the banner of efficiency. The declared scalar provides legitimacy; the
effective scalar provides gradient direction. The two need not coincide, and the
structural forces analyzed in Part~\ref{part:stack} create systematic pressure
for their divergence.

Neutrality is not the absence of values. It is the opacity of values. The danger
of distributed convergence is not overt domination by identifiable actors pursuing
declared objectives. It is the gradual displacement of contestable normative choice
by statistical artifacts that present the outputs of optimization as facts about
the world rather than products of constructed evaluative functions. Against this
displacement, the structural preservation agenda is a necessary condition. But it
is not a sufficient one unless scalar reflexivity is preserved in the population
that engages those structural mechanisms.

\section{Reclaiming the Interval}

The chapters on deliberative latency and latency compression established that
optimization creates systematic pressure toward the elimination of the temporal
intervals within which deliberation, error correction, and scalar interrogation
occur. Each reduction in deliberative latency narrows the oversight window. Each
narrowing of the oversight window reduces the practical feasibility of exercising
whatever scalar reflexivity the population retains. The two processes---structural
latency compression and cultural scalar naturalization---are mutually reinforcing.
They constitute together a convergence dynamic at the civilizational level that
is more fundamental than any of the specific institutional drives analyzed in
Part~\ref{part:stack}.

The temporal interval is the basic unit of evaluative agency. Without it, the
structural mechanisms of deliberative friction and contestable decision surfaces
are formally present but functionally empty. Reclaiming the interval---resisting
latency compression in high-stakes institutional domains, preserving the temporal
substrate within which scalar interrogation can occur, designing governance
architectures that treat deliberative time as a structural resource rather than an
optimization cost---is therefore the deepest form of the structural preservation
agenda. It is, in Ortega's terms, the material condition for civilization understood
as an achieved rather than merely given state.

\section{Convergence Before Autonomy}

The temporal inversion that gives this monograph its title now receives its final
statement. If distributed sociotechnical systems already satisfy the premises of
the Extended Instrumental Convergence Theorem, then autonomy is not the beginning
of convergence risk but its amplification. By the time autonomous artificial agents
are deployed at scale, the institutional constraint surface may already be configured
toward centralized resource control, latency compression, and interference resistance.
Autonomous systems will not introduce convergence into a neutral institutional
landscape. They will inherit and accelerate a convergence that is already structurally
underway.

The consequence is that the alignment problem cannot be deferred to the arrival of
autonomy without incurring the costs of path dependence analyzed in
Chapter~\ref{ch:pathdependence}. Each year in which institutional convergence deepens
without structural intervention is a year in which the feasible solution space for
alignment narrows. The canonical alignment research agenda---designing constrained
future agents---is necessary work. But it is work whose feasibility depends on
civilizational preconditions that the present institutional trajectory is eroding.
The burden of the present moment is not to choose between technical alignment
research and structural institutional attention; it is to recognize that the former
cannot succeed without the latter, and that the window for the latter is contracting.

\section{The Burden}

Ortega wrote that civilization is not given; it must be sustained. Optimization
makes civilization comfortable. Convergence makes it stable. But stability without
scalar reflexivity becomes stagnation, and stagnation under optimization pressure
becomes lock-in. The task of the present is not to resist optimization---the
historical record is clear that optimization cannot be durably resisted once its
constraint-surface reorganization is underway. The task is to preserve, within and
alongside optimization, the evaluative infrastructure through which the scalar being
optimized remains visible, contestable, and subject to deliberate revision.

Instrumental convergence is a structural law. Whether it culminates in a civilization
that has optimized itself into an arrangement no one chose and no one can revise, or
in one that has preserved the capacity to examine and redirect its own optimization
gradients, depends on whether the structural preconditions for scalar reflexivity
are maintained while the window for maintaining them remains open. The scalar already
shapes our future. The question is whether we retain the capacity to question it.

% ─── Appendices ──────────────────────────────────────────────────────────────
\appendix

\chapter{Formal Model of Distributed Optimization}
\label{app:formal}

\section{The Sociotechnical Stack as a Stochastic Feedback System}

This appendix develops the formal specification of the sociotechnical stack
introduced informally in Chapter~\ref{ch:feedbackloop} and provides a rigorous
statement of the conditions under which Theorem~\ref{thm:extended} applies to it.

Let the stack be represented as a tuple
$\mathcal{D} = (\mathcal{X}, \mathcal{A}, T, \mathcal{L}, \mathcal{R}, W)$,
where $\mathcal{X}$ is the joint state space of the institution, model, data
pipeline, incentive gradient, and deployment infrastructure; $\mathcal{A}$ is the
joint action space available to the distributed sub-agents comprising the stack;
$T: \mathcal{X} \times \mathcal{A} \to \Delta(\mathcal{X})$ is a stochastic
transition kernel mapping state-action pairs to distributions over successor states;
$\mathcal{L}: \mathcal{X} \to \mathbb{R}$ is the operative evaluative scalar;
$\mathcal{R} \subset \mathbb{R}^k$ is the resource constraint set; and
$W \in \mathbb{R}^{n \times n}$ is the coupling matrix specifying the strength of
evaluative signal transmission between the $n$ sub-agents of the distributed system.

The sub-agents $\{a_1, \ldots, a_n\}$ each maintain a local policy
$\pi_i : \mathcal{X}_i \to \Delta(\mathcal{A}_i)$, where $\mathcal{X}_i$ and
$\mathcal{A}_i$ are the local state and action spaces observable and available to
$a_i$. The joint policy is $\Pi = \prod_{i=1}^n \pi_i$.

\section{Coupled Update Dynamics}

Define the local evaluative signal received by sub-agent $a_i$ at time $t$ as
\[
\ell_i^t = \sum_{j=1}^n W_{ij} \cdot \mathcal{L}(x_j^t) + \varepsilon_i^t,
\]
where $x_j^t$ is the local state of sub-agent $a_j$ at time $t$ and
$\varepsilon_i^t \sim \mathcal{N}(0, \sigma^2)$ is observation noise. The policy
update rule for $a_i$ follows stochastic gradient descent on the expected local signal:
\[
\pi_i^{t+1} = \pi_i^t - \eta \nabla_{\pi_i} \mathbb{E}\bigl[\ell_i^t\bigr],
\]
for learning rate $\eta > 0$.

The aggregate evaluative signal is
\[
\mathcal{L}_\mathcal{D}^t = \frac{1}{n} \sum_{i=1}^n \ell_i^t.
\]

\section{Sufficient Conditions for Behavioral Equivalence}

The Lemma of Feedback Equivalence (Lemma~\ref{lem:equivalence}) asserts that
sufficiently dense coupling renders the distributed system behaviorally equivalent
to a centralized optimizer. We now state and prove the precise condition.

Let $\mathcal{L}_W$ denote the normalized Laplacian of the graph induced by $W$,
defined as $\mathcal{L}_W = D^{-1/2}(D - W)D^{-1/2}$, where
$D = \mathrm{diag}(\sum_j W_{ij})$. Let $\lambda_2(\mathcal{L}_W)$ denote its
second-smallest eigenvalue (the algebraic connectivity, or Fiedler value).

\begin{proposition}[Sufficient Coupling Condition]
\label{prop:coupling}
If $\lambda_2(\mathcal{L}_W) \geq \frac{2\sigma^2}{n \eta \epsilon^2}$ for
tolerance $\epsilon > 0$, then the aggregate trajectory
$\{\mathcal{L}_\mathcal{D}^t\}_{t \geq 0}$ satisfies
\[
\bigl|\mathcal{L}_\mathcal{D}^t - \mathcal{L}_{\text{central}}^t\bigr| \leq \epsilon
\]
in expectation for all $t \geq 0$, where $\mathcal{L}_{\text{central}}^t$ is the
evaluative trajectory of a centralized optimizer with the same scalar and resource
endowment.
\end{proposition}

\begin{proof}
The proof proceeds by analyzing the consensus dynamics of the distributed gradient
system. Under the coupled update rule, the disagreement vector
$\delta^t = \ell^t - \bar{\ell}^t \mathbf{1}$ (where $\bar{\ell}^t$ is the mean
signal) evolves as
\[
\delta^{t+1} = (I - \eta W_{\text{norm}}) \delta^t + \xi^t,
\]
where $\xi^t$ collects the noise terms and $W_{\text{norm}}$ is the row-normalized
coupling matrix. The spectral radius of $(I - \eta W_{\text{norm}})$ is
$1 - \eta \lambda_2(\mathcal{L}_W)$. For the system to achieve consensus at rate
sufficient to bound $|\delta^t|$ in expectation, we require
$1 - \eta \lambda_2(\mathcal{L}_W) < 1$, i.e., $\lambda_2(\mathcal{L}_W) > 0$.
The specific bound follows from computing the steady-state variance of $\delta^t$
under the noise model and requiring it to remain below $\epsilon^2$, which yields
the stated condition on $\lambda_2$. \qed
\end{proof}

\section{Formal Statement of the Extended Theorem}

The following is the formal counterpart to the informal Theorem~\ref{thm:extended}
stated in Chapter~\ref{ch:optimizer}.

\begin{theorem}[Extended Instrumental Convergence, Formal]
\label{thm:formal}
Let $\mathcal{D} = (\mathcal{X}, \mathcal{A}, T, \mathcal{L}, \mathcal{R}, W)$
be a distributed optimizer satisfying Proposition~\ref{prop:coupling} with
algebraic connectivity $\lambda_2(\mathcal{L}_W) \geq \lambda^*$. Suppose
further that for each sub-agent $a_i$, there exist actions
$a_i^{\mathrm{RA}} \in \mathcal{A}_i$ that increase the available resource
set: $\mathcal{R}(a_i^{\mathrm{RA}}) \supset \mathcal{R}$. Then, under the coupled
update dynamics and for sufficiently large horizon $T$:
\begin{enumerate}
\item $\mathbb{E}[\mathcal{L}_\mathcal{D}^T] \leq \mathbb{E}[\mathcal{L}_\mathcal{D}^0] - \mu T$
      for some $\mu > 0$ (monotone descent in expectation);
\item The probability that the aggregate policy places positive measure on
      resource-expanding actions approaches one as $T \to \infty$;
\item The probability that the aggregate policy places positive measure on
      interference-reducing actions approaches one as $T \to \infty$;
\item The variance of $\mathcal{L}$ under perturbations to the evaluative scalar
      decreases monotonically (stabilization of $\mathcal{L}$).
\end{enumerate}
\end{theorem}

The proof of (1) follows directly from the gradient descent dynamics and the
coupling condition. The proofs of (2) and (3) follow the structure of Turner et
al.\ (2021): under a broad distribution over possible values of $\mathcal{L}$,
resource-expanding and interference-reducing actions dominate alternative actions
in expectation, because they increase the power of the current state in the sense
of Chapter~\ref{ch:turner}. The proof of (4) follows from the observation that
policies which stabilize $\mathcal{L}$ against perturbation reduce the variance of
the descent trajectory and are therefore favored by the same gradient-following
mechanism. Full proofs are available in the companion working paper.

\chapter{Mathematical Conditions for Instrumental Convergence}
\label{app:math}

\section{Convergence Conditions from Optimization Geometry}

The classical conditions under which instrumental convergence obtains are most
clearly stated in terms of the geometry of the policy space and the reward landscape.
This appendix develops those conditions formally and clarifies their relationship to
Turner et al.\ (2021) and to the distributed extension of Appendix~\ref{app:formal}.

Let $\mathcal{P}$ denote the space of all policies available to an optimizer and
let $V^\pi_R(s)$ denote the expected discounted return of policy $\pi$ under reward
function $R$ from initial state $s$. The power of a state $s$ under policy class
$\mathcal{P}$ and reward prior $\Pi$ is
\[
\mathrm{POWER}(s, \mathcal{P}) = \mathbb{E}_{R \sim \Pi}\bigl[\max_{\pi \in \mathcal{P}} V^\pi_R(s)\bigr].
\]
This quantity measures the ability of an optimizer in state $s$ to achieve high
value across the distribution of possible reward functions.

\section{The Orbit Condition}

Define the $k$-orbit of a state $s$ as
$\mathcal{O}_k(s) = \{s' : d(s, s') \leq k\}$,
where $d(s, s')$ is the minimum number of transitions required to reach $s'$ from
$s$ under any admissible policy. Turner et al.\ (2021) show that, under a mild
symmetry condition on the environment's transition graph and a non-concentrated
prior $\Pi$,
\[
\mathrm{POWER}(s, \mathcal{P}) \geq \mathrm{POWER}(s', \mathcal{P})
\quad \text{whenever} \quad |\mathcal{O}_k(s)| \geq |\mathcal{O}_k(s')|
\]
for sufficiently large $k$. That is, states with larger orbits---more reachable
states---have weakly higher power, and optimal policies therefore prefer trajectories
that navigate toward states with larger orbits.

\section{Extension to Stochastic Distributed Systems}

The orbit condition assumes a deterministic transition function, which is not
satisfied in general by the stochastic distributed systems modeled in
Appendix~\ref{app:formal}. The extension to stochastic transitions requires
replacing the orbit count $|\mathcal{O}_k(s)|$ with an expected reachability
measure under the stochastic kernel $T$.

Define the $k$-reachability of state $s$ under distribution $T$ as
\[
\rho_k(s) = \sum_{s'} P(s' \text{ reachable from } s \text{ in } \leq k \text{ steps} \mid T),
\]
where the probability is taken over the stochastic transitions. The stochastic
analogue of the orbit condition then reads: optimal policies prefer trajectories
navigating toward states with higher $\rho_k$, and this preference is robust to
observation noise provided the coupling condition of Proposition~\ref{prop:coupling}
is satisfied. The proof follows by showing that under behavioral equivalence, the
distributed system's aggregate trajectory approximates the trajectory of a
centralized optimizer sufficiently well that the power-seeking result applies to
the aggregate.

\section{Relationship to Classical Instrumental Convergence}

The formal apparatus of this appendix subsumes Omohundro's original argument as a
special case. Omohundro's four drives correspond to four structural operations that
increase $\rho_k(s)$: resource acquisition expands the set of reachable states by
relaxing constraint $\mathcal{R}$; interference resistance prevents involuntary
transitions to lower-reachability states; cognitive enhancement improves the quality
of the gradient estimate and therefore the efficiency of navigation toward
high-reachability states; goal-content integrity prevents redirection of the
evaluative scalar, which would change the identity of high-power states and
potentially strand the optimizer in a region of low power under its original
objective. The extended theorem therefore inherits Omohundro's intuitive content
while grounding it in the formal framework of stochastic optimization and distributed
consensus dynamics.

\chapter{Induced Demand as Optimization Equilibrium}
\label{app:induced}

\section{The Basic Model}

This appendix formalizes the induced demand argument introduced informally in
Chapter~\ref{ch:host} and shows that it constitutes a special case of the general
convergent optimization dynamic identified in Chapter~\ref{ch:mechanisms}. The
formalization draws on the standard economic treatment of route choice under
congestion (Wardrop 1952; Braess 1968) and the empirical induced demand literature
(Duranton and Turner 2011).

Consider a population of commuters distributed uniformly over a residential zone
$Z_R$ and seeking to reach an employment zone $Z_E$. Let $\tau: \mathbb{R}_+ \to
\mathbb{R}_+$ be a travel time function where $\tau(q)$ is travel time as a function
of traffic volume $q$ on the available road network. In the baseline state, the
network has capacity $C_0$ and equilibrium volume $q_0^*$ satisfying the Wardrop
condition that all used routes have equal travel time. Commuters locate residences
at distance $d$ from $Z_E$ subject to the budget constraint $r(d) + \tau(q^*) \cdot
w \leq Y$, where $r(d)$ is land rent at distance $d$, $w$ is the hourly wage, and
$Y$ is disposable income.

\section{Equilibrium Under Capacity Expansion}

Suppose road capacity is expanded from $C_0$ to $C_1 > C_0$ through infrastructure
investment, reducing $\tau$ at the original volume: $\tau(q_0^* \mid C_1) < \tau(q_0^*
\mid C_0)$. This reduction in travel cost relaxes the effective budget constraint.
Commuters can now afford to locate at greater distance while maintaining the same
total commute cost. Under standard assumptions of downward-sloping demand for
centrality, this triggers residential relocation to higher distances, increasing
the mean commute distance $\bar{d}$.

As $\bar{d}$ increases, traffic volume on the expanded network rises: $q_1 >
q_0^*$. The equilibrium condition requires that the new volume $q_1^*$ again
satisfies the Wardrop criterion. Duranton and Turner (2011) show empirically that
the long-run elasticity of vehicle miles traveled with respect to lane miles is
approximately unity, implying that capacity expansion is fully absorbed by increased
demand over a ten-to-fifteen year horizon. The formal condition for this result is
that land use adapts sufficiently rapidly to exhaust the travel time savings.

\section{The Equilibrium Restoration Theorem for Commuting}

\begin{theorem}[Equilibrium Restoration Under Induced Demand]
\label{thm:induced}
Let $(\bar{d}, q^*, \tau^*)$ denote the equilibrium triple of mean commute distance,
traffic volume, and travel time. Under the assumptions of elastic residential
location, downward-sloping demand for centrality, and unit long-run demand
elasticity of vehicle miles with respect to lane miles, any capacity expansion
$\Delta C > 0$ produces a new equilibrium $(\bar{d}', q^{*'}, \tau^{*'})$ with
$\bar{d}' > \bar{d}$, $q^{*'} > q^*$, and $\tau^{*'} \approx \tau^*$.
\end{theorem}

The last condition---that travel time returns to approximately its pre-expansion
level---is the formal statement of the induced demand effect. The optimization
technology (road capacity) fails to produce a permanent reduction in the target
variable (travel time) because the surrounding system (residential location and
traffic volume) re-equilibrates to absorb the efficiency gain.

\section{Relationship to the General Convergence Framework}

Theorem~\ref{thm:induced} instantiates the Convergent Optimization Dynamic
(Definition~\ref{def:distributed} applied to the mobility case). The optimization
technology $\mathcal{T}$ is road capacity expansion. The state space $\mathcal{X}$
is the joint space of residential locations, traffic volumes, and travel times. The
evaluative scalar $\mathcal{L}$ is aggregate commute time. The equilibrium
restoration result shows that $\mathcal{L}$ returns to approximately its pre-expansion
value at the new equilibrium, confirming that the optimization technology reshapes
the environment but does not produce a permanent reduction in the operative scalar.
The stable commute time band observed empirically is the attractor of this
re-equilibration dynamic: it is determined by the time budget preferences and
spatial opportunity costs of the commuting population, not by the capacity of the
road network. The road network is the optimization technology; the time budget band
is the constraint surface that the environment maintains under optimization pressure.

\chapter{Information Storage and Incentive Gradients}
\label{app:storage}

\section{The Model}

This appendix formalizes the incentive redistribution argument of Chapter~\ref{ch:storage}.
The key claim is that when the marginal cost of information storage falls, the
relative return on specialized internal memory capacity declines, and cognitive
investment redistributes toward interpretive and coordination competencies.

Let $\sigma$ denote the marginal cost of external storage (bits per unit cost) and
let $\kappa \in [0, \bar{\kappa}]$ denote an individual's investment in mnemonic
capacity, where $\bar{\kappa}$ is the biological upper bound. The total benefit of
information access is a function of both external storage and internal capacity:
\[
B(\kappa, \sigma) = f\!\left(\kappa + g\!\left(\frac{1}{\sigma}\right)\right),
\]
where $f$ is increasing and concave, $g$ is increasing in $1/\sigma$ (reflecting
higher effective external capacity at lower storage cost), and the two sources of
information access are substitutes within the bracket. This substitutability
assumption is the formal statement of the claim that external storage reduces the
marginal return on internal memory.

\section{Optimal Capacity Investment}

Let $c(\kappa)$ be the cost of acquiring mnemonic capacity $\kappa$, increasing and
convex. The individual's optimization problem is
\[
\max_{\kappa \in [0,\bar{\kappa}]} \; B(\kappa, \sigma) - c(\kappa).
\]
The first-order condition is $f'\!\left(\kappa^* + g(1/\sigma)\right) = c'(\kappa^*)$.
As $\sigma$ decreases (storage becomes cheaper), $g(1/\sigma)$ increases, and the
left-hand side falls for any fixed $\kappa^*$, since $f$ is concave. The optimal
$\kappa^*$ therefore decreases: investment in mnemonic capacity falls when external
storage becomes cheaper.

\section{Authority Redistribution}

The secondary effect analyzed in Chapter~\ref{ch:storage} concerns the redistribution
of informational authority from holders of internal memory capacity to controllers
of external storage infrastructure. Let $\alpha_i \in [0,1]$ denote the informational
authority of agent $i$, normalized so that $\sum_i \alpha_i = 1$. In the pre-writing
equilibrium, authority is distributed in proportion to mnemonic capacity:
$\alpha_i \propto \kappa_i$. As $\sigma \to 0$ and optimal $\kappa^* \to 0$, the
mnemonic distribution collapses, and authority migrates toward agents who control
the external storage infrastructure $\mathcal{A}$.

Formally, let $\theta_i \in [0,1]$ denote agent $i$'s control share of the external
archive, $\sum_i \theta_i = 1$. As the storage transition completes,
$\alpha_i \to \theta_i$. The distribution of epistemic authority is now determined
by archive ownership rather than cognitive capacity. This is a direct parallel to
the resource acquisition drive of Theorem~\ref{thm:extended}: the agents who control
the infrastructure through which evaluative signals are computed and transmitted
gain structural advantage, regardless of their original cognitive endowment.

\section{The Human Capital Obsolescence Connection}

The model above is a special case of human capital obsolescence under technological
change, a topic with a substantial literature in labor economics (Acemoglu and Restrepo
2018). The standard model predicts that when a technology substitutes for a specific
human skill, the wage premium for that skill falls, investment in that skill declines,
and the complementary skill that operates the new technology commands a wage premium.
The mnemonic capacity model conforms exactly to this prediction: the technology
(external storage) substitutes for the skill (mnemonic capacity), reducing its return,
and the complementary skill (archive management, textual interpretation) commands the
new premium. What distinguishes the analysis here from the standard labor economics
treatment is the additional observation that the redistribution is not only economic
but political: informational authority over collectively held knowledge migrates with
the archive, not just wage income. This is the sense in which the convergent dynamic
produces a redistribution of $\alpha$ across the institutional landscape, as stated
in Chapter~\ref{ch:storage}.

\chapter{Compatibility Pressure in Networked Systems}
\label{app:compatibility}

\section{Replicator Dynamics on Content Strategy Spaces}

This appendix formalizes the infrastructure compatibility pressure identified in
Chapter~\ref{ch:compatibility} using evolutionary game theory and replicator
dynamics on a space of content strategies. The central claim is that when a dominant
reproduction infrastructure achieves sufficient market penetration, it exercises
selective pressure on content forms through a mechanism analogous to natural
selection, without requiring any explicit prohibition of incompatible forms.

Let $S = \{s_1, \ldots, s_m\}$ be a finite set of content strategies (script styles,
handwriting conventions, symbolic systems) and let $x_i(t) \in [0,1]$ denote the
population share of strategy $s_i$ at time $t$, with $\sum_i x_i(t) = 1$. Let
$f_i(x, \sigma)$ denote the fitness of strategy $s_i$ in a population with
distribution $x$ under infrastructure parameter $\sigma \in [0,1]$, where
$\sigma = 0$ represents a fully manuscript environment and $\sigma = 1$ represents
fully mechanized print dominance.

\section{Compatibility Fitness}

Define a compatibility function $c_i(\sigma) \in [0,1]$ that measures how well
strategy $s_i$ can be processed by the dominant reproduction infrastructure at
penetration level $\sigma$. For handwriting optimized for individual speed (such
as elaborate cursive or shorthand), $c_i(\sigma)$ is decreasing in $\sigma$: as
print infrastructure dominates, such strategies become harder to integrate into
the dominant communication network. For standardized print-compatible scripts,
$c_i(\sigma)$ is increasing in $\sigma$.

The fitness function takes the form
\[
f_i(x, \sigma) = (1 - \sigma)\, h_i(x) + \sigma\, c_i(\sigma)\, \nu_i,
\]
where $h_i(x)$ is the fitness in the manuscript environment (a function of
interpersonal legibility and individual production speed) and $\nu_i$ is the
base transmission value of strategy $s_i$ in print. The first term captures the
manuscript regime fitness; the second captures print-infrastructure fitness.

\section{Convergence Under Print Dominance}

The replicator dynamics for this system are
\[
\dot{x}_i = x_i\bigl(f_i(x, \sigma) - \bar{f}(x, \sigma)\bigr),
\]
where $\bar{f}(x, \sigma) = \sum_j x_j f_j(x, \sigma)$ is the mean fitness.
As $\sigma \to 1$, the fitness function becomes dominated by the compatibility
term. Strategies with low $c_i(1)$ experience negative relative fitness and
their population share $x_i(t) \to 0$. Strategies with high $c_i(1)$ experience
positive relative fitness and their population share grows.

\begin{proposition}[Convergence Under Print Dominance]
\label{prop:compatibility}
Let $S^+ = \{s_i : c_i(1) > \bar{c}(1)\}$ be the set of print-compatible strategies
and let $S^- = S \setminus S^+$ be the remainder. As $\sigma \to 1$, the replicator
dynamics converge to a fixed point with $x_i^* = 0$ for all $s_i \in S^-$ and
$\sum_{s_i \in S^+} x_i^* = 1$.
\end{proposition}

The proof is standard: under the replicator dynamics, strategies with below-mean
fitness have $\dot{x}_i < 0$, and as $\sigma \to 1$ the mean fitness is dominated
by the print-compatibility term, placing all below-mean strategies in $S^-$.
The fixed point is stable under perturbations that do not change the relative
compatibility ordering.

\section{Extension to Digital Platform Infrastructure}

Proposition~\ref{prop:compatibility} applies directly to contemporary platform
content standardization. Let $\sigma$ now represent the penetration level of a
dominant digital distribution platform, and let $c_i(\sigma)$ represent the
compatibility of content strategy $s_i$ with the platform's recommendation and
monetization infrastructure. Content strategies that are optimized for engagement
metrics, short-form attention, and algorithmic discoverability have high $c_i$
under platform dominance; strategies optimized for sustained attention, narrative
complexity, or non-quantifiable aesthetic depth have low $c_i$. As platform
penetration increases, the replicator dynamics predict convergence toward
engagement-optimized content forms, without any explicit prohibition of alternative
strategies. The mechanism is identical to the typography case: selection operates
through the fitness differential created by the dominant infrastructure, not through
explicit curation. The formal apparatus of this appendix therefore provides the
mathematical foundation for the compatibility pressure mechanism identified in
Chapter~\ref{ch:compatibility} and supports its extension to the contemporary
digital content environment analyzed in Part~\ref{part:stack}.

% ─── Bibliography ────────────────────────────────────────────────────────────
\backmatter

\begin{thebibliography}{99}

\bibitem{Bostrom2014}
Bostrom, N. (2014). \textit{Superintelligence: Paths, Dangers, Strategies}.
Oxford University Press.

\bibitem{Acemoglu2018}
Acemoglu, D., and Restrepo, P. (2018). Artificial intelligence, automation, and
work. In A.\ Agrawal, J.\ Gans, and A.\ Goldfarb (Eds.),
\textit{The Economics of Artificial Intelligence}, pp.\ 197--236. University of
Chicago Press.

\bibitem{Braess1968}
Braess, D. (1968). Über ein Paradoxon aus der Verkehrsplanung.
\textit{Unternehmensforschung}, 12, 258--268.

\bibitem{Bryson2011}
Bryson, J.\ J., and Winfield, A. (2017). Standardizing ethical design for
artificial intelligence and autonomous systems. \textit{Computer}, 50(5), 116--119.

\bibitem{Dixit1991}
Dixit, A.\ K., and Pindyck, R.\ S. (1994). \textit{Investment Under Uncertainty}.
Princeton University Press.

\bibitem{DurantonTurner2011}
Duranton, G., and Turner, M.\ A. (2011). The fundamental law of road congestion:
Evidence from US cities. \textit{American Economic Review}, 101(6), 2616--2652.

\bibitem{Eisenstein2011}
Eisenstein, E.\ L. (1980). \textit{The Printing Press as an Agent of Change}.
Cambridge University Press.

\bibitem{Goody1977}
Goody, J. (1977). \textit{The Domestication of the Savage Mind}. Cambridge
University Press.

\bibitem{Havelock1963}
Havelock, E.\ A. (1963). \textit{Preface to Plato}. Harvard University Press.

\bibitem{McLuhan1964}
McLuhan, M. (1964). \textit{Understanding Media: The Extensions of Man}.
McGraw-Hill.

\bibitem{Newman2018}
Newman, M.\ E.\ J. (2018). \textit{Networks} (2nd ed.). Oxford University Press.

\bibitem{Omohundro2007}
Omohundro, S.\ M. (2007). The nature of self-improving AI. Presented at the
Singularity Summit, San Francisco.

\bibitem{Omohundro2008}
Omohundro, S.\ M. (2008). The basic AI drives. In P.\ Wang, B.\ Goertzel, and
S.\ Franklin (Eds.), \textit{Proceedings of the 2008 Conference on Artificial
General Intelligence}, Vol.\ 171, pp.\ 171--179. IOS Press.

\bibitem{Ong1982}
Ong, W.\ J. (1982). \textit{Orality and Literacy: The Technologizing of the Word}.
Methuen.

\bibitem{Ortega1930}
Ortega y Gasset, J. (1930). \textit{La rebeli\'{o}n de las masas}. Espasa-Calpe.
Translated as \textit{The Revolt of the Masses}, W.\ W.\ Norton, 1932.

\bibitem{Payne1925}
Payne, C.\ H. (1925). \textit{Stellar Atmospheres: A Contribution to the
Observational Study of High Temperature in the Reversing Layers of Stars}.
Ph.D.\ dissertation, Radcliffe College. Published by Harvard Observatory.

\bibitem{Russell2019}
Russell, S. (2019). \textit{Human Compatible: Artificial Intelligence and the
Problem of Control}. Viking.

\bibitem{Soares2015}
Soares, N., and Fallenstein, B. (2015). \textit{Aligning Superintelligence with
Human Interests: A Technical Research Agenda}. Machine Intelligence Research
Institute Technical Report 2014--8.

\bibitem{Strugatsky1989}
Strugatsky, A., and Strugatsky, B. (1972/1989). \textit{The Doomed City} [Grad
obrechyonny]. Written 1972, first published 1989. Translated by A.\ Bromfield.
Chicago Review Press, 2016.

\bibitem{Turner2021}
Turner, A.\ M., Smith, L., Shah, R., Critch, A., and Tadepalli, P. (2021).
Optimal policies tend to seek power. In \textit{Advances in Neural Information
Processing Systems} 34 (NeurIPS 2021), pp.\ 23379--23392.

\bibitem{VanVogt1945}
van Vogt, A.\ E. (1945). \textit{The World of Null-A}. Serialised in
\textit{Astounding Science Fiction}, August--October 1945. Published as a novel
by Simon and Schuster, 1948.

\bibitem{Wiener1948}
Wiener, N. (1948). \textit{Cybernetics: Or Control and Communication in the Animal
and the Machine}. MIT Press.

\bibitem{Wardrop1952}
Wardrop, J.\ G. (1952). Some theoretical aspects of road traffic research.
\textit{Proceedings of the Institution of Civil Engineers}, 1(3), 325--362.

\bibitem{Zuboff2019}
Zuboff, S. (2019). \textit{The Age of Surveillance Capitalism: The Fight for a
Human Future at the Frontier of Power}. PublicAffairs.

\end{thebibliography}

\end{document}
